\documentclass[aps,prd,twocolumn,superscriptaddress,nofootinbib]{revtex4-2}

\usepackage{amsmath}
\usepackage{mathtools}
\usepackage{amssymb}
\usepackage{amsthm}
\usepackage{booktabs}
\usepackage{tabularx}
\usepackage{xcolor}
\usepackage{graphicx}
\usepackage{mathrsfs}

\newtheorem{theorem}{Theorem}[section]
\newtheorem{proposition}[theorem]{Proposition}
\newtheorem{lemma}[theorem]{Lemma}
\newtheorem{corollary}[theorem]{Corollary}
\newtheorem{conjecture}[theorem]{Conjecture}
\theoremstyle{definition}
\newtheorem{definition}[theorem]{Definition}
\theoremstyle{remark}
\newtheorem{remark}[theorem]{Remark}
\newtheorem{example}[theorem]{Example}

\usepackage{hyperref}
\hypersetup{
    pdftitle={Arithmetische Positivitaet und absolute Hodge-Klassen: Warum Hodge-Riemann-Positivitaet Delignes Absolutheit nicht verstaerkt},
    pdfauthor={Lukas Geiger},
    colorlinks=true,
    linkcolor=black,
    urlcolor=blue,
    citecolor=black
}

\begin{document}

% ===================================================================
% TITELSEITE
% ===================================================================

\title{Arithmetische Positivitaet und absolute Hodge-Klassen:\\Warum Hodge-Riemann-Positivitaet\\Delignes Absolutheit nicht verstaerkt}

\author{Lukas Geiger}
\email{https://orcid.org/0009-0005-7296-1534}
\affiliation{Unabhaengiger Forscher, Bernau im Schwarzwald, Deutschland}

\date{\today}

\begin{abstract}
Wir fuehren das Konzept der \textit{arithmetischen Positivitaet} fuer die Hodge-Vermutung ein: eine rationale $(p,p)$-Klasse $\alpha \in H^{2p}(X, \mathbb{Q})$ auf einer glatten projektiven Varietaet $X/\mathbb{C}$ heisst arithmetisch positiv, wenn sie gleichzeitig Hodge-Riemann-Positivitaet unter allen Polarisierungen, Delignes Absolutheit unter allen Automorphismen von $\mathbb{C}$ und Lefschetz-Kompatibilitaet erfuellt. Wir beweisen, dass algebraische Klassen arithmetisch positiv sind, und etablieren Funktorialitaet unter endlichen Morphismen und K\"unneth-Produkten primitiver Klassen. Unser \textbf{Hauptresultat ist negativ}: der arithmetisch positive Kegel stimmt exakt mit Delignes absoluten Hodge-Klassen ueberein ($\mathrm{AP}^p(X) = \mathrm{AbsHodge}^p(X)$), weil die Hodge-Riemann-Bilinearrelationen die Positivitaetsbedingungen fuer jede Klasse vom Typ $(p,p)$ automatisch garantieren. Der natuerliche Versuch, Delignes Absolutheit durch Hodge-Riemann-Positivitaet zu verstaerken, liefert daher keine neuen Einschraenkungen. Die Arithmetische Positivitaetsvermutung reduziert sich auf Delignes Frage (1982): ,,Sind alle absoluten Hodge-Klassen algebraisch?'' Ferner sind die Obstruktionen (O1) und (O3) in der No-Go-Formulierung fuer rationale $(p,p)$-Klassen vakuos; nur (O2), Versagen der Absolutheit, ist effektiv. Wir fuehren das Galois-Hodge-Riemann-Spektrum als diagnostisches Werkzeug ein und diskutieren abelsche Varietaeten sowie Richtungen zur Suche nach genuein staerkeren Positivitaetsbedingungen jenseits der Hodge-Riemann-Relationen.
\end{abstract}

\keywords{Hodge-Vermutung, algebraische Zykel, arithmetische Positivitaet, absolute Hodge-Klassen, Hodge-Riemann-Bilinearform, No-Go-Theorem}

\thanks{KI-Werkzeuge (Claude, Anthropic) wurden fuer mathematische Formalisierung und Texterstellung verwendet. Alle mathematischen Inhalte, Vermutungen und die Verantwortung fuer die Korrektheit liegen ausschliesslich beim Autor.}

\maketitle


% ===================================================================
% 1. EINFUEHRUNG
% ===================================================================
\section{Einfuehrung}
\label{sec:intro}

\subsection{Die Hodge-Vermutung}

Sei $X$ eine glatte projektive Varietaet der Dimension $n$ ueber $\mathbb{C}$. Die Hodge-Zerlegung
\begin{equation}
H^k(X, \mathbb{C}) = \bigoplus_{p+q=k} H^{p,q}(X)
\label{eq:hodge_decomp}
\end{equation}
drueckt die $k$-te singulaere Kohomologie als direkte Summe von Dolbeault-Kohomologiegruppen aus. Eine Klasse $\alpha \in H^{2p}(X, \mathbb{Q})$ heisst \textit{rationale Hodge-Klasse}, wenn ihr Bild in $H^{2p}(X, \mathbb{C})$ vollstaendig in $H^{p,p}(X)$ liegt. Jeder algebraische Zykus $Z$ der Kodimension $p$ definiert eine rationale Hodge-Klasse $\mathrm{cl}(Z) \in H^{2p}(X, \mathbb{Q}) \cap H^{p,p}(X)$ ueber die Zykusklassenabbildung.

\begin{conjecture}[Hodge, 1950]
\label{conj:hodge}
Jede rationale Hodge-Klasse auf einer glatten projektiven Varietaet ueber $\mathbb{C}$ ist eine $\mathbb{Q}$-Linearkombination von Klassen algebraischer Zykel.
\end{conjecture}

Trotz ueber sieben Jahrzehnten der Forschung bleibt die Vermutung im Allgemeinen offen. Sie ist bekannt fuer:
\begin{itemize}
\item Divisoren ($p = 1$): durch das Lefschetz-$(1,1)$-Theorem.
\item Abelsche Varietaeten der Dimension $\leq 5$: durch Arbeiten von Hazama, Moonen und anderen.
\item Produkte elliptischer Kurven und bestimmte Fermat-Hyperflaechen.
\item Uniruled-Varietaeten (fuer Grad-$2$-Klassen).
\end{itemize}

\subsection{Warum Konstruktion scheitert}

Alle grossen Ansaetze zur Hodge-Vermutung teilen eine gemeinsame Strategie: den algebraischen Zykus, der eine gegebene Hodge-Klasse repraesentiert, zu \textit{konstruieren}. Diese Arbeit argumentiert, dass diese Strategie grundlegend fehlgeleitet ist, und schlaegt eine Alternative vor: zu beweisen, dass nichtalgebraische Hodge-Klassen ueber eine Positivitaetsobstruktion \textit{unmoeglich} sind.

Die strukturellen Versagen bestehender Ansaetze sind:

\textbf{1. Direkte Algebraisierung.} Gegeben eine rationale $(p,p)$-Klasse, versuche eine Untervarietaet zu konstruieren, deren Fundamentalklasse mit ihr uebereinstimmt. Dies scheitert, weil es keine kanonische Projektion von Hodge-Kohomologie auf algebraische Zykel gibt und die Konstruktion unter Deformation nicht funktoriell ist.

\textbf{2. Variationen von Hodge-Strukturen.} Verwende Familien $X_t$ und verfolge, wie sich Hodge-Klassen aendern. Dies scheitert, weil Hodge-Klassen ,,springen'' koennen -- die Bedingung, vom Typ $(p,p)$ zu sein, ist nicht offen in der Zariski-Topologie -- und Monodromie naive Stabilitaetsargumente zerstoert \cite{Voisin2002}.

\textbf{3. Motivische Methoden.} Ersetze Varietaeten durch Motive und hoffe, dass Hodge-Klassen motivisch sind. Dies scheitert, weil die Kategorie der Motive nicht vollstaendig konstruiert ist: Grothendiecks Standardvermutungen \cite{Grothendieck1969} bleiben offen, und die Beziehung zwischen Hodge-Klassen und absoluten Hodge-Klassen ist ungeklaert \cite{Andre1996}.

\textbf{4. Reduktion mod $p$.} Reduziere auf Charakteristik $p$, berufe die Tate-Vermutung und lifte zurueck. Dies scheitert, weil die Bruecke Hodge $\Rightarrow$ Tate $\Rightarrow$ Hodge nicht geschlossen ist: Frobenius-Invarianz impliziert nicht den Hodge-Typ, und das Liften algebraischer Zykel ist nicht eindeutig \cite{Charles2013}.

Wie wir unten zeigen, stoesst auch die in dieser Arbeit vorgeschlagene Positivitaetsalternative auf eine fundamentale Einschraenkung: die natuerlichen Positivitaetsbedingungen reduzieren sich auf Delignes Absolutheit. Dennoch markiert die Analyse eine praezise Grenze fuer diese Klasse von Ansaetzen.

\subsection{Die Positivitaetsalternative}

Wir beobachten, dass grosse Durchbrueche in benachbarten Vermutungen einem anderen Muster folgen:

\begin{center}
\begin{tabular}{ll}
\hline
Vermutung & Schluesseltechnik \\
\hline
Weil-Vermutungen & Positivitaet der Frobenius-Eigenwerte \\
BSD (Rang 1) & Positivitaet der Hoehen (Gross--Zagier) \\
Tate (abelsche Var.) & Positivitaet via Faltings' Theorem \\
\hline
\end{tabular}
\end{center}

Jeder dieser Erfolge beruht nicht auf der Konstruktion des gewuenschten Objekts, sondern darauf zu zeigen, dass seine Nichtexistenz eine Positivitaetsbedingung verletzen wuerde. Der Hodge-Vermutung fehlt genau dies: \textit{eine Positivitaetsstruktur, die Algebraizitaet erzwingt, ohne Zykel zu konstruieren}.

Diese Arbeit fuehrt eine solche Struktur ein -- \textit{arithmetische Positivitaet} -- die Hodge-Riemann-Positivitaet unter allen Polarisierungen mit Delignes Absolutheit unter allen Automorphismen von $\mathbb{C}$ kombiniert. Ein zentrales Ergebnis ist jedoch \textit{negativ}: die Hodge-Riemann-Bilinearrelationen garantieren die Positivitaetsbedingungen fuer jede Klasse vom Typ $(p,p)$ automatisch, sodass der arithmetisch positive Kegel mit Delignes absoluten Hodge-Klassen uebereinstimmt: $\mathrm{AP}^p(X) = \mathrm{AbsHodge}^p(X)$. Der Versuch, Absolutheit durch Hodge-Riemann-Positivitaet zu verstaerken, liefert keine neuen Einschraenkungen. Die Arithmetische Positivitaetsvermutung reduziert sich auf Delignes Frage: ,,Sind alle absoluten Hodge-Klassen algebraisch?''

Wir betrachten dieses negative Resultat als informativ: es markiert eine praezise Grenze fuer positivitaetsbasierte Ansaetze zur Hodge-Vermutung und identifiziert das Galois-Hodge-Riemann-Spektrum als diagnostisches Werkzeug.

\subsection{Gliederung}

Abschnitt~\ref{sec:background} wiederholt die Hodge-Riemann-Bilinearform und absolute Hodge-Klassen. Abschnitt~\ref{sec:arithmetic_positivity} fuehrt arithmetische Positivitaet ein. Abschnitt~\ref{sec:easy_direction} beweist, dass algebraische Klassen arithmetisch positiv sind. Abschnitt~\ref{sec:functoriality} etabliert Funktorialitaet. Abschnitt~\ref{sec:nogo} formuliert das No-Go-Theorem. Abschnitt~\ref{sec:abelian} beweist die Aequivalenz fuer abelsche Varietaeten. Abschnitt~\ref{sec:obstruction} identifiziert die praezise Obstruktion. Abschnitt~\ref{sec:discussion} diskutiert die verbleibende Luecke und die Beziehung zu Delignes Programm.


% ===================================================================
% 2. HINTERGRUND
% ===================================================================
\section{Hintergrund}
\label{sec:background}

\subsection{Hodge-Riemann-Bilinearform}

Sei $X$ eine glatte projektive Varietaet der Dimension $n$, ausgestattet mit einer K\"ahler-Form $\omega$, die eine ample Klasse $L = [\omega] \in H^{1,1}(X) \cap H^2(X, \mathbb{Z})$ definiert. Der \textit{Lefschetz-Operator} $\Lambda_L: H^k(X) \to H^{k+2}(X)$ ist definiert durch $\Lambda_L(\alpha) = L \cup \alpha$.\footnote{Wir verwenden die Konvention $\Lambda_L(\alpha) = L \cup \alpha$; einige Autoren (z.\,B.\ Voisin~\cite{Voisin2002}, Griffiths--Harris~\cite{GH1978}) reservieren $\Lambda$ fuer den dualen Lefschetz-Operator (den Adjungierten von $L$ bezueglich des Hodge-Skalarprodukts) und notieren die Cup-Produkt-Operation einfach als $L$ oder $L \wedge$.}

\begin{theorem}[Hard Lefschetz]
\label{thm:hard_lefschetz}
Fuer $k \leq n$ ist die Abbildung $\Lambda_L^{n-k}: H^k(X, \mathbb{Q}) \to H^{2n-k}(X, \mathbb{Q})$ ein Isomorphismus.
\end{theorem}

Eine Klasse $\alpha \in H^k(X, \mathbb{Q})$ mit $k \leq n$ ist \textit{primitiv} (bezueglich $L$), wenn $\Lambda_L^{n-k+1}(\alpha) = 0$. Die \textit{Hodge-Riemann-Bilinearform} auf der primitiven Kohomologie $P^k(X, \mathbb{Q}) \subset H^k(X, \mathbb{Q})$ ist:
\begin{equation}
Q_L(\alpha, \beta) = \int_X \alpha \cup \beta \cup L^{n-k}\,.
\label{eq:HR_form}
\end{equation}

\begin{theorem}[Hodge-Riemann-Bilinearrelationen]
\label{thm:HR}
Auf dem primitiven Unterraum $P^{p,q}(X) = P^{p+q}(X) \cap H^{p,q}(X)$ mit $p + q = k$ ist die hermitesche Form
\begin{equation}
h_L(\alpha, \beta) = i^{p-q}\,(-1)^{k(k-1)/2}\, Q_L(\alpha, \bar{\beta})
\label{eq:HR_hermitian}
\end{equation}
positiv definit.
\end{theorem}

Die Hodge-Riemann-Relationen sind das fundamentale Positivitaetsresultat in der Hodge-Theorie. Sie liefern Definitheit \textit{relativ zu einer festen Polarisierung} $L$. Ein zentrales Thema dieser Arbeit ist, dass Algebraizitaet Positivitaet nicht nur fuer ein $L$ erfordert, sondern fuer \textit{alle} zulaessigen Polarisierungen und Operationen.

\subsection{Absolute Hodge-Klassen}

Deligne \cite{Deligne1982} fuehrte eine Verfeinerung von Hodge-Klassen ein. Ein Automorphismus $\sigma \in \mathrm{Aut}(\mathbb{C})$ sendet eine glatte projektive Varietaet $X/\mathbb{C}$ auf die konjugierte Varietaet $X^\sigma$ (gleiche Gleichungen, Koeffizienten verdrillt durch $\sigma$). Der Vergleichsisomorphismus liefert:
\begin{equation}
H^{2p}_{\mathrm{dR}}(X/\mathbb{C}) \xrightarrow{\sim} H^{2p}_{\mathrm{dR}}(X^\sigma/\mathbb{C})\,,
\end{equation}
der die de~Rham-Klasse von $\alpha$ auf $X$ auf eine Klasse $\alpha^\sigma$ auf $X^\sigma$ abbildet.

\begin{definition}[Deligne]
\label{def:absolute}
Eine rationale Hodge-Klasse $\alpha \in H^{2p}(X, \mathbb{Q}) \cap H^{p,p}(X)$ ist \textit{absolut} (oder \textit{absolut Hodge}), wenn fuer jeden $\sigma \in \mathrm{Aut}(\mathbb{C})$ die konjugierte Klasse $\alpha^\sigma$ wieder eine Hodge-Klasse auf $X^\sigma$ ist.
\end{definition}

\begin{theorem}[Deligne \cite{Deligne1982}]
\label{thm:deligne_absolute}
\begin{enumerate}
\item Jede algebraische Klasse ist absolut.
\item Auf abelschen Varietaeten ist jede Hodge-Klasse absolut.
\end{enumerate}
\end{theorem}

Delignes Resultat zeigt, dass Absolutheit eine \textit{notwendige} Bedingung fuer Algebraizitaet ist. Die Frage ist, ob sie hinreichend ist. Unser Konzept der arithmetischen Positivitaet \textit{versucht}, Absolutheit durch Hinzufuegen von Hodge-Riemann-Positivitaet unter allen Konjugationen zu verstaerken; wie in Abschnitt~\ref{sec:discussion} gezeigt, gelingt dies nicht (die zusaetzlichen Bedingungen sind automatisch), aber die Analyse ist informativ.


% ===================================================================
% 3. ARITHMETISCHE POSITIVITAET
% ===================================================================
\section{Arithmetische Positivitaet}
\label{sec:arithmetic_positivity}

Wir fuehren nun das zentrale Konzept ein.

\subsection{Definition}

Sei $X$ eine glatte projektive Varietaet der Dimension $n$ ueber $\mathbb{C}$, und sei $\mathrm{Amp}(X)$ der ample Kegel von $X$.

\begin{definition}[Arithmetische Positivitaet]
\label{def:arith_pos}
Eine rationale Hodge-Klasse $\alpha \in H^{2p}(X, \mathbb{Q}) \cap H^{p,p}(X)$ ist \textit{arithmetisch positiv}, wenn sie die folgenden drei Bedingungen gleichzeitig erfuellt:

\begin{enumerate}
\item[\textbf{(AP1)}] \textbf{Universelle Hodge-Riemann-Positivitaet.} Fuer jede ample Klasse $L \in \mathrm{Amp}(X)$ erfuellt die primitive Komponente $\alpha_0$ von $\alpha$ (bezueglich $L$) die \textit{vorzeichenbehaftete} Hodge-Riemann-Ungleichung:
\begin{equation}
(-1)^p\, Q_L(\alpha_0, \alpha_0) \geq 0\,,
\label{eq:AP1}
\end{equation}
wobei $Q_L$ die Hodge-Riemann-Form~\eqref{eq:HR_form} ist. Das Vorzeichen $(-1)^p$ ist wesentlich: die Hodge-Riemann-Bilinearrelationen (Theorem~\ref{thm:HR}) ergeben $(-1)^p Q_L > 0$ auf primitiven $(p,p)$-Klassen, nicht $Q_L > 0$.

\item[\textbf{(AP2)}] \textbf{Absolute Stabilitaet.} Fuer jeden $\sigma \in \mathrm{Aut}(\mathbb{C})$:
\begin{enumerate}
\item[(a)] $\alpha^\sigma$ ist eine Hodge-Klasse auf $X^\sigma$ (d.h., $\alpha$ ist absolut).
\item[(b)] Die konjugierte Klasse $\alpha^\sigma$ erfuellt (AP1) auf $X^\sigma$ bezueglich jeder amplen Klasse $L^\sigma \in \mathrm{Amp}(X^\sigma)$.
\end{enumerate}

\item[\textbf{(AP3)}] \textbf{Lefschetz-Kompatibilitaet.} Fuer alle $0 \leq j \leq n - 2p$ und alle amplen $L$, sei $\gamma_j = \Lambda_L^j \alpha \in H^{2(p+j)}(X)$ und $\gamma_{j,0}$ ihre primitive Komponente bezueglich $L$ in $H^{2(p+j)}(X)$. Dann:
\begin{equation}
(-1)^{p+j}\, Q_{L}(\gamma_{j,0},\, \gamma_{j,0}) \geq 0\,.
\label{eq:AP3}
\end{equation}
\end{enumerate}

Die Menge der arithmetisch positiven Klassen in $H^{2p}(X, \mathbb{Q}) \cap H^{p,p}(X)$ wird mit $\mathrm{AP}^p(X)$ bezeichnet.
\end{definition}

\begin{remark}
\label{rem:redundancy}
Die Bedingungen (AP1), (AP2b) und (AP3) sind \textit{nicht} unabhaengig von (AP2a). Wie in Proposition~\ref{prop:absolute_implies_ap} unten gezeigt, garantieren die Hodge-Riemann-Bilinearrelationen (AP1), (AP2b) und (AP3) automatisch fuer jede Klasse vom Typ $(p,p)$. Die einzige substantielle Bedingung ist daher (AP2a): Delignes Absolutheit. Insbesondere gilt $\mathrm{AP}^p(X) = \mathrm{AbsHodge}^p(X)$ (Bemerkung~\ref{rem:ap_equals_absolute}). Die Definition behaelt alle drei Bedingungen aus darstellerischen Gruenden bei: sie machen den \textit{Typ} der Positivitaet explizit, den algebraische Klassen geniessen, auch wenn die Hodge-Riemann-Relationen sie fuer $(p,p)$-Klassen automatisch machen. Genuein staerkere Bedingungen zu finden -- jenseits der Hodge-Riemann-Relationen -- die \textit{tatsaechlich} ueber Absolutheit hinaus einschraenken, bleibt ein offenes Problem (siehe Abschnitt~\ref{sec:discussion}).
\end{remark}

\subsection{Die Arithmetische Positivitaetsvermutung}

\begin{conjecture}[Arithmetische Positivitaetsvermutung]
\label{conj:AP}
Sei $X$ eine glatte projektive Varietaet ueber $\mathbb{C}$. Dann gilt:
\begin{equation}
\mathrm{AP}^p(X) = \mathrm{cl}\!\left(\mathrm{CH}^p(X)_\mathbb{Q}\right)\,,
\end{equation}
wobei $\mathrm{CH}^p(X)_\mathbb{Q}$ die Chow-Gruppe der Kodimension-$p$-Zykel mit rationalen Koeffizienten ist und $\mathrm{cl}$ die Zykusklassenabbildung.
\end{conjecture}

\begin{proposition}
\label{prop:equivalence}
Die folgenden Beziehungen gelten zwischen der APV und der HV:
\begin{enumerate}
\item Die Hodge-Vermutung impliziert die Arithmetische Positivitaetsvermutung.
\item Die ,,schwere Richtung'' der APV ($\mathrm{AP}^p(X) \subseteq \mathrm{cl}(\mathrm{CH}^p(X)_\mathbb{Q})$) zusammen mit der Aussage, dass alle rationalen Hodge-Klassen arithmetisch positiv sind ($H^{2p}(X, \mathbb{Q}) \cap H^{p,p}(X) \subseteq \mathrm{AP}^p(X)$), impliziert die HV.
\item Unbedingt: $\mathrm{cl}(\mathrm{CH}^p(X)_\mathbb{Q}) \subseteq \mathrm{AP}^p(X)$ (Theorem~\ref{thm:easy_direction} unten).
\end{enumerate}
\end{proposition}

\begin{proof}
(1) Wenn die HV gilt, ist jede Hodge-Klasse algebraisch, und algebraische Klassen erfuellen (AP1)--(AP3) (Theorem~\ref{thm:easy_direction} unten). Somit $\mathrm{AP}^p(X) \subseteq H^{2p}(X, \mathbb{Q}) \cap H^{p,p}(X) = \mathrm{cl}(\mathrm{CH}^p(X)_\mathbb{Q})$, und die umgekehrte Inklusion gilt nach Theorem~\ref{thm:easy_direction}, also $\mathrm{AP}^p(X) = \mathrm{cl}(\mathrm{CH}^p(X)_\mathbb{Q})$.

(2) Dies erfordert zwei separate Annahmen. \textit{Erstens}, die schwere Richtung der APV gelte: $\mathrm{AP}^p(X) \subseteq \mathrm{cl}(\mathrm{CH}^p(X)_\mathbb{Q})$, d.h.\ jede arithmetisch positive Klasse ist algebraisch. \textit{Zweitens}, jede rationale Hodge-Klasse sei arithmetisch positiv: $H^{2p}(X, \mathbb{Q}) \cap H^{p,p}(X) \subseteq \mathrm{AP}^p(X)$. Unter beiden Annahmen liegt jede rationale Hodge-Klasse $\alpha$ in $\mathrm{AP}^p(X)$ (nach der zweiten Annahme) und damit in $\mathrm{cl}(\mathrm{CH}^p(X)_\mathbb{Q})$ (nach der ersten Annahme), also ist $\alpha$ algebraisch. Dies ergibt die HV.

Keine der beiden Annahmen genuegt allein: die APV ($\mathrm{AP}^p = \mathrm{alg}^p$) zeigt, dass AP-Klassen algebraisch sind, aber ohne das Wissen, dass alle Hodge-Klassen AP sind, kann man nicht schliessen, dass alle Hodge-Klassen algebraisch sind.

(3) Siehe Theorem~\ref{thm:easy_direction}.
\end{proof}

\begin{remark}
Die APV allein impliziert die HV nicht direkt. Die APV behauptet $\mathrm{AP}^p = \mathrm{alg}^p$, waehrend die HV $\mathrm{Hodge}^p = \mathrm{alg}^p$ behauptet. Da $\mathrm{AP}^p \subseteq \mathrm{Hodge}^p$ (arithmetische Positivitaet stellt zusaetzliche Bedingungen jenseits einer Hodge-Klasse), gibt die APV $\mathrm{alg}^p \subseteq \mathrm{Hodge}^p$ (trivial), aber nicht $\mathrm{Hodge}^p \subseteq \mathrm{alg}^p$ ohne das zusaetzliche Wissen, dass alle Hodge-Klassen arithmetisch positiv sind. Der nichttriviale Inhalt zerfaellt in zwei Teile: (i)~zeige dass nicht-AP Hodge-Klassen nicht existieren (d.h. jede Hodge-Klasse ist AP), und (ii)~zeige dass AP Algebraizitaet impliziert (die schwere Richtung der APV).
\end{remark}

Der nichttriviale Inhalt der APV ist die ,,schwere Richtung'': $\mathrm{AP}^p(X) \subseteq \mathrm{cl}(\mathrm{CH}^p(X)_\mathbb{Q})$. Diese besagt, dass arithmetische Positivitaet Algebraizitaet \textit{erzwingt}.


% ===================================================================
% 4. DIE LEICHTE RICHTUNG
% ===================================================================
\section{Algebraische Klassen sind arithmetisch positiv}
\label{sec:easy_direction}

\begin{theorem}[Leichte Richtung]
\label{thm:easy_direction}
Sei $X$ eine glatte projektive Varietaet ueber $\mathbb{C}$. Jede Klasse im Bild der Zykusklassenabbildung ist arithmetisch positiv:
\begin{equation}
\mathrm{cl}\!\left(\mathrm{CH}^p(X)_\mathbb{Q}\right) \subseteq \mathrm{AP}^p(X)\,.
\end{equation}
\end{theorem}

\begin{proof}
Sei $\alpha = \mathrm{cl}(Z)$ fuer einen algebraischen Zykus $Z$ der Kodimension $p$ mit rationalen Koeffizienten. Wir verifizieren (AP1)--(AP3).

\textit{(AP1):} Sei $\alpha_0$ die $L$-primitive Komponente von $\alpha = \mathrm{cl}(Z)$. Da $\alpha$ algebraisch ist, liegt es in $H^{p,p}(X) \cap H^{2p}(X, \mathbb{Q})$, und seine primitive Komponente $\alpha_0$ liegt in $P^{p,p}(X) \cap H^{2p}(X, \mathbb{Q})$. Nach den Hodge-Riemann-Bilinearrelationen (Theorem~\ref{thm:HR}) ist die hermitesche Form $h_L$ positiv definit auf $P^{p,p}(X)$. Fuer eine rationale $(p,p)$-Klasse gilt $\bar{\alpha}_0 = \alpha_0$, $i^{p-p} = 1$ und $(-1)^{k(k-1)/2} = (-1)^{p(2p-1)}$ mit $k = 2p$. Also $h_L(\alpha_0, \alpha_0) = (-1)^{p(2p-1)}\, Q_L(\alpha_0, \alpha_0)$. Da $(-1)^{p(2p-1)} = (-1)^{2p^2 - p} = (-1)^p$ (weil $(-1)^{2p^2} = 1$), ergeben die Hodge-Riemann-Relationen:
\begin{equation}
(-1)^p\, Q_L(\alpha_0, \alpha_0) > 0
\end{equation}
falls $\alpha_0 \neq 0$. Fuer $\alpha_0 = 0$ (d.h., wenn $\alpha$ keine primitive Komponente hat) gilt die Ungleichung trivial. Dieses Argument verwendet nur die Hodge-Riemann-Relationen auf $P^{p,p}(X)$ und erfordert nicht, dass $\alpha_0$ selbst die Klasse eines effektiven Zykus ist.

\textit{(AP2):} Teil (a) ist Delignes Theorem (Theorem~\ref{thm:deligne_absolute}): algebraische Klassen sind absolut. Teil (b) folgt, weil $\sigma$ algebraische Zykel auf $X$ auf algebraische Zykel auf $X^\sigma$ sendet (Algebraizitaet ist eine Eigenschaft polynomialer Gleichungen, die kovariant unter $\sigma$ transformieren), und algebraische Klassen auf $X^\sigma$ (AP1) durch das obige Argument erfuellen.

\textit{(AP3):} Das Hard-Lefschetz-Theorem ist kompatibel mit der Zykusklassenabbildung: wenn $\alpha = \mathrm{cl}(Z)$, dann ist $\Lambda_L^j \alpha = \mathrm{cl}(Z \cdot L^j)$, wobei $Z \cdot L^j$ der Schnitt mit $j$ Kopien eines Hyperebenenschnitts ist. Dies ist wieder ein algebraischer Zykus der Kodimension $p + j$. Seine $L$-primitive Komponente $\gamma_{j,0}$ liegt in $P^{p+j,\,p+j}(X)$. Da $\gamma_{j,0}$ der primitive Anteil einer algebraischen Klasse ist, ergeben die Hodge-Riemann-Relationen $(-1)^{p+j} Q_L(\gamma_{j,0}, \gamma_{j,0}) \geq 0$, mit demselben Argument wie in (AP1).
\end{proof}


% ===================================================================
% 5. FUNKTORIALITAET
% ===================================================================
\section{Funktorialitaet der Arithmetischen Positivitaet}
\label{sec:functoriality}

Damit der No-Go-Ansatz funktioniert, muss der arithmetisch positive Kegel funktoriell sein: er muss unter den Standardoperationen der algebraischen Geometrie erhalten bleiben.

\begin{proposition}[Rueckzug -- endliche Morphismen]
\label{prop:pullback}
Sei $f: Y \to X$ ein endlicher Morphismus glatter projektiver Varietaeten. Wenn $\alpha \in \mathrm{AP}^p(X)$, dann gilt $f^*\alpha \in \mathrm{AP}^p(Y)$.
\end{proposition}

\begin{proof}
Sei $f: Y \to X$ endlich vom Grad $d$.

\textit{(AP1):} Da $f$ endlich ist, ist $f^*L_X$ ampel auf $Y$ fuer jede ample Klasse $L_X \in \mathrm{Amp}(X)$. Fuer die primitive Komponente $\alpha_0$ von $\alpha$ bezueglich $L_X$ liefert die Projektionsformel:
\[
\int_Y f^*\alpha_0 \cup f^*\bar{\alpha}_0 \cup (f^*L_X)^{n-2p} = d \int_X \alpha_0 \cup \bar{\alpha}_0 \cup L_X^{n-2p},
\]
wobei $n = \dim X = \dim Y$. Da $d > 0$ und die rechte Seite $(-1)^p Q_{L_X}(\alpha_0, \alpha_0) \geq 0$ nach Voraussetzung erfuellt, erhalten wir $(-1)^p Q_{f^*L_X}(f^*\alpha_0, f^*\alpha_0) \geq 0$. Fuer eine allgemeine ample Klasse $L_Y \in \mathrm{Amp}(Y)$ ist die Menge $\{L \in \mathrm{Amp}(Y) : (-1)^p Q_L(\beta_0, \beta_0) \geq 0\}$ (wobei $\beta_0$ die $L$-primitive Komponente von $f^*\alpha$ ist) abgeschlossen im amplen Kegel (Stetigkeit von $Q_L$ in $L$; vgl.\ Voisin~\cite{Voisin2002}, Prop.\ 6.32). Da $f$ endlich ist, bildet $f^*$ ampel auf ampel ab, und $\{f^*L_X : L_X \in \mathrm{Amp}(X)\}$ ist ein Unterkegel von $\mathrm{Amp}(Y)$. Die Erweiterung auf ganz $\mathrm{Amp}(Y)$ folgt durch die Dichte von $f^*\mathrm{Amp}(X) + \mathrm{Amp}(Y)$ und die Abgeschlossenheit der Positivitaetsbedingung.

\textit{(AP2):} Rueckzug kommutiert mit $\sigma$-Konjugation: $(f^*\alpha)^\sigma = (f^\sigma)^*(\alpha^\sigma)$, und $f^\sigma: Y^\sigma \to X^\sigma$ ist wieder endlich vom selben Grad. Da $\alpha^\sigma \in \mathrm{AP}^p(X^\sigma)$ nach Voraussetzung, folgt das Resultat aus dem (AP1)-Argument angewandt auf $f^\sigma$ und $\alpha^\sigma$.

\textit{(AP3):} Der Lefschetz-Operator erfuellt $\Lambda_{f^*L}^j \circ f^* = f^* \circ \Lambda_L^j$ (Kompatibilitaet von Rueckzug mit Cup-Produkt). Daher gilt $\Lambda_{f^*L}^j(f^*\alpha) = f^*(\Lambda_L^j \alpha)$, und die Positivitaet der primitiven Komponente folgt aus dem Grad-$d$-Argument angewandt auf $\Lambda_L^j\alpha$.
\end{proof}

\begin{remark}
Die im Beweis verwendete Gradformel $Q_{f^*L_X}(f^*\alpha_0, f^*\alpha_0) = \deg(f) \cdot Q_{L_X}(\alpha_0, \alpha_0)$ setzt voraus, dass $f$ ein endlicher Morphismus ist (aequivalent: generisch endlich mit $\dim Y = \dim X$). Fuer einen generisch endlichen Morphismus $f: Y \to X$ mit $\dim Y = \dim X$, der nicht endlich ist, liefert die Projektionsformel die Gradrelation zwar auf einer dichten offenen Teilmenge, aber der Verzweigungsort erfordert zusaetzliche Sorgfalt; die Aussage von Proposition~\ref{prop:pullback} ist daher auf endliche Morphismen beschraenkt.
\end{remark}

\begin{remark}
Fuer Morphismen $f: Y \to X$ mit $\dim Y > \dim X$ ist der Rueckzug $f^*L_X$ nicht ampel auf $Y$, und die primitive Zerlegung auf $Y$ erfordert eine relative Hodge-Riemann-Analyse. Die Funktorialitaet von $\mathrm{AP}^p$ unter allgemeinen Morphismen bleibt offen.
\end{remark}

\begin{proposition}[Vorwaertsschub -- endliche Morphismen]
\label{prop:pushforward}
Sei $f: Y \to X$ ein endlicher Morphismus glatter projektiver Varietaeten (also $\dim Y = \dim X$). Wenn $\alpha \in \mathrm{AP}^{p}(Y)$, dann gilt $f_*\alpha \in \mathrm{AP}^p(X)$.
\end{proposition}

\begin{proof}[Beweisskizze]
Die Projektionsformel $f_*f^* = \deg(f) \cdot \mathrm{id}$ auf $H^*(X, \mathbb{Q})$ impliziert, dass $f_*$ auf dem Bild von $f^*$ injektiv ist (bis auf Skalar). Fuer $\alpha \in \mathrm{AP}^p(Y)$ und jede ample Klasse $L_X \in \mathrm{Amp}(X)$ ist $f^*L_X$ ampel auf $Y$ (da $f$ endlich). Die Hodge-Riemann-Form auf $X$ ist mit jener auf $Y$ durch die Adjunktion $\int_X f_*\beta \cup \gamma = \int_Y \beta \cup f^*\gamma$ verknuepft. Zusammen mit der Gradformel folgen die Vorzeichenbedingungen fuer $f_*\alpha$ aus jenen von $\alpha$. Die absolute Stabilitaet uebertraegt sich, da $f^\sigma$ wieder endlich vom selben Grad ist.
\end{proof}

\begin{remark}
Fuer Morphismen $f: Y \to X$ mit $\dim Y > \dim X$ erfordert der Vorwaertsschub $f_*$ ein Gysin-Argument, das die Hodge-Riemann-Form auf $Y$ mit jener auf $X$ ueber die Grothendieck-Riemann-Roch-Formel und die relative primitive Zerlegung in Beziehung setzt. Diese Analyse wird hier nicht durchgefuehrt; die Funktorialitaet von $\mathrm{AP}^p$ unter allgemeinem Vorwaertsschub bleibt offen.
\end{remark}

\begin{proposition}[K\"unneth-Produkt -- primitive Klassen]
\label{prop:kunneth}
Seien $\alpha \in \mathrm{AP}^p(X)$ und $\beta \in \mathrm{AP}^q(Y)$ primitiv bezueglich amplen Klassen $L_X$ bzw.\ $L_Y$. Dann gilt $\alpha \boxtimes \beta \in \mathrm{AP}^{p+q}(X \times Y)$.
\end{proposition}

\begin{remark}
Die Erweiterung auf nichtprimitive Klassen erfordert eine sorgfaeltige Analyse der Kreuzterme in der Lefschetz-Zerlegung auf dem Produkt, die hier nicht durchgefuehrt wird. Das Resultat fuer primitive Klassen deckt den strukturell wesentlichen Fall ab.
\end{remark}

Wir etablieren zunaechst ein Schluessellemma ueber die Primitivitaet von aeusseren Produkten.

\begin{lemma}[Primitivitaet aeusserer Produkte]
\label{lem:primitive_product}
Seien $X$ und $Y$ glatte projektive Varietaeten der Dimensionen $n$ bzw.\ $m$, mit amplen Klassen $L_X \in \mathrm{Amp}(X)$ und $L_Y \in \mathrm{Amp}(Y)$. Setze $L = L_X \boxtimes 1 + 1 \boxtimes L_Y \in \mathrm{Amp}(X \times Y)$. Wenn $\alpha \in P^{2p}(X, \mathbb{Q})$ primitiv bezueglich $L_X$ und $\beta \in P^{2q}(Y, \mathbb{Q})$ primitiv bezueglich $L_Y$ ist, dann ist $\alpha \boxtimes \beta$ primitiv bezueglich $L$ auf $X \times Y$.
\end{lemma}

\begin{proof}
Setze $k = p + q$ und $N = n + m = \dim(X \times Y)$. Wir muessen zeigen, dass $L^{N-2k+1} \cup (\alpha \boxtimes \beta) = 0$. Da $L = L_X \boxtimes 1 + 1 \boxtimes L_Y$, liefert die Binomialentwicklung fuer beliebiges $j \geq 0$:
\[
L^j \cup (\alpha \boxtimes \beta) = \sum_{a+b=j} \binom{j}{a} (L_X^a \cup \alpha) \boxtimes (L_Y^b \cup \beta).
\]
Jeder Summand verschwindet, es sei denn beide Faktoren sind von null verschieden. Da $\alpha$ $L_X$-primitiv ist, gilt $L_X^{n-2p+1} \cup \alpha = 0$, also $L_X^a \cup \alpha = 0$ fuer $a \geq n - 2p + 1$. Analog gilt $L_Y^b \cup \beta = 0$ fuer $b \geq m - 2q + 1$. Ein Summand mit $a + b = j$ ist nur dann von null verschieden, wenn $a \leq n - 2p$ und $b \leq m - 2q$, was $j \leq (n - 2p) + (m - 2q) = N - 2k$ erfordert. Daher gilt $L^j \cup (\alpha \boxtimes \beta) = 0$ fuer alle $j \geq N - 2k + 1$, was genau die Primitivitaetsbedingung ist.
\end{proof}

\begin{lemma}[Vorzeichenkompatibilitaet fuer aeussere Produkte]
\label{lem:sign_product}
Unter den Voraussetzungen von Lemma~\ref{lem:primitive_product}, mit $k = p + q$:
\begin{equation}
(-1)^k\, Q_L(\alpha \boxtimes \beta,\, \alpha \boxtimes \beta) = \binom{N-2k}{n-2p}\, \bigl[(-1)^p\, Q_{L_X}(\alpha, \alpha)\bigr] \cdot \bigl[(-1)^q\, Q_{L_Y}(\beta, \beta)\bigr],
\label{eq:sign_kunneth}
\end{equation}
wobei $N = n + m$ und $Q_L$, $Q_{L_X}$, $Q_{L_Y}$ die Hodge-Riemann-Formen~\eqref{eq:HR_form} auf $X \times Y$, $X$ bzw.\ $Y$ sind.
\end{lemma}

\begin{proof}
Nach Lemma~\ref{lem:primitive_product} ist $\alpha \boxtimes \beta$ $L$-primitiv, also liefert die Definition~\eqref{eq:HR_form} direkt:
\[
Q_L(\alpha \boxtimes \beta,\, \alpha \boxtimes \beta) = \int_{X \times Y} (\alpha \boxtimes \beta)^2 \cup L^{N-2k}.
\]
Entwicklung von $L^{N-2k}$ ueber die Binomialformel wie in Lemma~\ref{lem:primitive_product} zeigt, dass der einzige nichtverschwindende Term bei $a = n - 2p$, $b = m - 2q$ auftritt (erzwungen durch $a \leq n-2p$, $b \leq m-2q$ und $a + b = N - 2k = (n-2p) + (m-2q)$). Nach der K\"unneth-Formel fuer Integration:
\[
\int_{X \times Y} (\alpha \boxtimes \beta)^2 \cup L^{N-2k} = \binom{N-2k}{n-2p} \left(\int_X \alpha^2 \cup L_X^{n-2p}\right) \left(\int_Y \beta^2 \cup L_Y^{m-2q}\right).
\]
Nach Definition~\eqref{eq:HR_form} sind die Integrale auf $X$ und $Y$ gerade $Q_{L_X}(\alpha, \alpha)$ bzw.\ $Q_{L_Y}(\beta, \beta)$. Daher:
\[
Q_L(\alpha \boxtimes \beta,\, \alpha \boxtimes \beta) = \binom{N-2k}{n-2p}\, Q_{L_X}(\alpha, \alpha) \cdot Q_{L_Y}(\beta, \beta).
\]
Multiplikation beider Seiten mit $(-1)^k = (-1)^{p+q} = (-1)^p \cdot (-1)^q$ ergibt Gleichung~\eqref{eq:sign_kunneth}.
\end{proof}

\begin{proof}[Beweis von Proposition~\ref{prop:kunneth}]
Wir verifizieren (AP1)--(AP3) fuer $\alpha \boxtimes \beta$ mit $\alpha \in \mathrm{AP}^p(X)$ primitiv und $\beta \in \mathrm{AP}^q(Y)$ primitiv.

\textit{(AP1):} Sei $L = L_X \boxtimes 1 + 1 \boxtimes L_Y$ fuer ample Klassen $L_X \in \mathrm{Amp}(X)$, $L_Y \in \mathrm{Amp}(Y)$. Nach Lemma~\ref{lem:primitive_product} ist $\alpha \boxtimes \beta$ $L$-primitiv, also gleicht es seiner eigenen primitiven Komponente. Nach Lemma~\ref{lem:sign_product}:
\[
(-1)^{p+q}\, Q_L(\alpha \boxtimes \beta,\, \alpha \boxtimes \beta) = \binom{N-2k}{n-2p}\, \bigl[(-1)^p\, Q_{L_X}(\alpha, \alpha)\bigr] \cdot \bigl[(-1)^q\, Q_{L_Y}(\beta, \beta)\bigr].
\]
Der Binomialkoeffizient $\binom{N-2k}{n-2p} > 0$, und beide geklammerten Faktoren sind $\geq 0$ durch (AP1) fuer $\alpha$ und $\beta$. Daher gilt $(-1)^{p+q}\, Q_L(\alpha \boxtimes \beta, \alpha \boxtimes \beta) \geq 0$. Ample Klassen der Form $L_X \boxtimes 1 + 1 \boxtimes L_Y$ bilden einen Unterkegel von $\mathrm{Amp}(X \times Y)$. Da die Menge $\{L \in \overline{\mathrm{Amp}}(X \times Y) : (-1)^{p+q} Q_L(\gamma_0, \gamma_0) \geq 0\}$ abgeschlossen ist ($Q_L$ haengt stetig von $L$ ab) und diesen Unterkegel enthaelt, erstreckt sich die Ungleichung auf den Abschluss. Fuer Produktvarietaeten, deren amplen Kegel durch Rueckzugsklassen erzeugt wird (z.B.\ wenn $X$ oder $Y$ Picard-Zahl~$1$ hat), genuegt dies. Im Allgemeinen erfordert die Erweiterung auf ganz $\mathrm{Amp}(X \times Y)$ die Verifikation der Ungleichung fuer Polarisierungen mit gemischten Termen, was offen bleibt.

\textit{(AP2):} Fuer beliebiges $\sigma \in \mathrm{Aut}(\mathbb{C})$ gilt $(\alpha \boxtimes \beta)^\sigma = \alpha^\sigma \boxtimes \beta^\sigma$ und $(X \times Y)^\sigma = X^\sigma \times Y^\sigma$. Da $\alpha^\sigma \in \mathrm{AP}^p(X^\sigma)$ und $\beta^\sigma \in \mathrm{AP}^q(Y^\sigma)$ nach Voraussetzung, laesst sich das (AP1)-Argument auf $\alpha^\sigma \boxtimes \beta^\sigma$ auf $X^\sigma \times Y^\sigma$ anwenden.

\textit{(AP3):} Der Lefschetz-Operator auf dem Produkt erfuellt $\Lambda_L^j(\alpha \boxtimes \beta) = \sum_{a+b=j} \binom{j}{a} (\Lambda_{L_X}^a \alpha) \boxtimes (\Lambda_{L_Y}^b \beta)$. Fuer primitive $\alpha$ und $\beta$ haben die nichttrivialen Terme $a \leq n-2p$ und $b \leq m-2q$, und die Positivitaet jedes solchen Terms folgt aus (AP3) fuer $\alpha$ und $\beta$ zusammen mit der Vorzeichenberechnung von Lemma~\ref{lem:sign_product}.

Fuer nichtprimitive Klassen fuehrt die primitive Zerlegung auf $X \times Y$ bezueglich $L$ Kreuzterme zwischen verschiedenen Lefschetz-Niveaus ein. Die vollstaendige Analyse dieser Kreuzterme bleibt offen; der obige Beweis deckt den strukturell wesentlichen Fall primitiver Klassen ab.
\end{proof}


% ===================================================================
% 6. DAS NO-GO-THEOREM
% ===================================================================
\section{Das No-Go-Theorem}
\label{sec:nogo}

Wir formulieren nun das zentrale Resultat: die Hodge-Vermutung, umformuliert als Positivitaetsobstruktion.

\begin{theorem}[Positivitaetsobstruktion -- Bedingt]
\label{thm:nogo}
Es gelte die Arithmetische Positivitaetsvermutung~\ref{conj:AP}. Sei $X$ eine glatte projektive Varietaet ueber $\mathbb{C}$, und sei $\alpha \in H^{2p}(X, \mathbb{Q}) \cap H^{p,p}(X)$ eine rationale Hodge-Klasse, die \textbf{keine} $\mathbb{Q}$-Linearkombination algebraischer Zykel ist. Dann existiert mindestens eine der folgenden Obstruktionen:

\begin{enumerate}
\item[\textbf{(O1)}] \textbf{Polarisierungsinstabilitaet:} Es existiert eine ample Klasse $L \in \mathrm{Amp}(X)$, sodass die primitive Komponente $\alpha_0$ von $\alpha$ die Bedingung $(-1)^p\,Q_L(\alpha_0, \alpha_0) < 0$ erfuellt.

\item[\textbf{(O2)}] \textbf{Galois-Instabilitaet:} Es existiert $\sigma \in \mathrm{Aut}(\mathbb{C})$, sodass entweder $\alpha^\sigma$ keine Hodge-Klasse auf $X^\sigma$ ist, oder $\alpha^\sigma$ (AP1) auf $X^\sigma$ verletzt.

\item[\textbf{(O3)}] \textbf{Lefschetz-Instabilitaet:} Es existieren $j \geq 1$ und $L \in \mathrm{Amp}(X)$, sodass $(-1)^{p+j}\,Q_L(\gamma_{j,0}, \gamma_{j,0}) < 0$, wobei $\gamma_{j,0}$ die $L$-primitive Komponente von $\Lambda_L^j \alpha \in H^{2(p+j)}(X)$ ist.
\end{enumerate}
\end{theorem}

\begin{proof}
Durch Kontraposition: wenn $\alpha$ (AP1)--(AP3) erfuellt, dann gilt $\alpha \in \mathrm{AP}^p(X) = \mathrm{cl}(\mathrm{CH}^p(X)_\mathbb{Q})$ durch die APV, im Widerspruch zur Annahme, dass $\alpha$ nicht algebraisch ist.
\end{proof}

\begin{remark}[Struktur der Obstruktion]
\label{rem:obstruction_structure}
Die drei Obstruktionen muessen im Licht der Hodge-Riemann-Relationen interpretiert werden.

\textbf{(O1) und (O3) sind fuer $(p,p)$-Klassen vakuos.} Die Hodge-Riemann-Bilinearrelationen (Theorem~\ref{thm:HR}) garantieren $(-1)^p Q_L(\alpha_0, \alpha_0) \geq 0$ fuer \textit{jede} rationale $(p,p)$-Klasse $\alpha$, nicht nur fuer algebraische. Die primitive Komponente $\alpha_0$ liegt in $P^{p,p}(X)$, und die Hodge-Riemann-Form ist dort positiv definit (mit dem Vorzeichen $(-1)^p$). Analog liegt die primitive Komponente $\gamma_{j,0}$ von $\Lambda_L^j \alpha$ in $P^{p+j,p+j}(X)$, sodass $(-1)^{p+j} Q_L(\gamma_{j,0}, \gamma_{j,0}) \geq 0$ automatisch gilt. Daher \textit{koennen Obstruktionen (O1) und (O3) fuer rationale Hodge-Klassen niemals eintreten}.

\textbf{(O2) ist die einzige effektive Obstruktion.} Der einzige Weg, auf dem eine Hodge-Klasse arithmetische Positivitaet verletzen kann, ist ueber (O2a): Versagen der Absolutheit ($\alpha^\sigma$ ist keine Hodge-Klasse auf $X^\sigma$ fuer ein $\sigma$). Teil (O2b) -- Versagen der Hodge-Riemann-Positivitaet auf $X^\sigma$ -- ist ebenfalls vakuos, da $\alpha^\sigma \in H^{p,p}(X^\sigma)$ die Hodge-Riemann-Ungleichungen automatisch erfuellt.

Folglich reduziert sich das No-Go-Theorem auf Delignes Beobachtung: nichtabsolute Hodge-Klassen (falls sie existieren) koennen nicht algebraisch sein.
\end{remark}

\subsection{Die unbedingte Aussage}

Waehrend Theorem~\ref{thm:nogo} von der APV abhaengt, koennen wir ein unbedingtes Teilresultat formulieren:

\begin{proposition}[Unbedingte Obstruktion]
\label{prop:unconditional}
Sei $\alpha \in H^{2p}(X, \mathbb{Q}) \cap H^{p,p}(X)$ eine rationale Hodge-Klasse. Wenn $\alpha$ (AP2a) nicht erfuellt -- d.h., es existiert $\sigma \in \mathrm{Aut}(\mathbb{C})$ mit $\alpha^\sigma \notin H^{p,p}(X^\sigma)$ -- dann ist $\alpha$ nicht algebraisch.
\end{proposition}

\begin{proof}
Dies ist Delignes Resultat (Theorem~\ref{thm:deligne_absolute}, Teil~1, Kontraposition): algebraische Klassen sind absolut.
\end{proof}

Wie in Abschnitt~\ref{sec:discussion} (Bemerkung~\ref{rem:ap_equals_absolute}) gezeigt, werden die Hodge-Riemann-Positivitaetsbedingungen (AP1), (AP2b) und (AP3) von jeder absoluten Hodge-Klasse automatisch erfuellt. Der Inhalt der APV jenseits von Proposition~\ref{prop:unconditional} ist daher identisch mit Delignes Frage: \textit{Sind alle absoluten Hodge-Klassen algebraisch?} Die ,,Luecke'' zwischen absoluten Hodge-Klassen und algebraischen Klassen wird durch die Positivitaetsbedingungen dieser Arbeit nicht verkleinert.


% ===================================================================
% 7. DER FALL ABELSCHER VARIETAETEN
% ===================================================================
\section{Abelsche Varietaeten}
\label{sec:abelian}

Abelsche Varietaeten bieten die guenstigste Umgebung fuer den arithmetischen Positivitaetsansatz, weil Delignes Theorem (alle Hodge-Klassen sind absolut) die Obstruktion (O2a) eliminiert.

\begin{theorem}
\label{thm:abelian}
Fuer abelsche Varietaeten ist die Arithmetische Positivitaetsvermutung~\ref{conj:AP} aequivalent zur Hodge-Vermutung. Genauer: HV fuer abelsche Varietaeten $\Rightarrow$ APV fuer abelsche Varietaeten (unbedingt nach Theorem~\ref{thm:easy_direction}), und APV fuer abelsche Varietaeten $\Rightarrow$ HV fuer abelsche Varietaeten (bedingt darauf, dass jede Hodge-Klasse auf einer abelschen Varietaet arithmetisch positiv ist; siehe Bemerkung unten).
\end{theorem}

\begin{proof}[Beweisstrategie (Skizze)]
Sei $A$ eine abelsche Varietaet der Dimension $g$, und sei $\alpha \in H^{2p}(A, \mathbb{Q}) \cap H^{p,p}(A)$.

\textit{Schritt 1:} Nach Delignes Theorem (Theorem~\ref{thm:deligne_absolute}, Teil~2) ist $\alpha$ absolut. Damit ist (AP2a) automatisch erfuellt.

\textit{Schritt 2:} Fuer abelsche Varietaeten hat die Hodge-Riemann-Form eine explizite Beschreibung in Termen der Riemannschen Form der Polarisierung. Sei $E$ die Riemannsche Form, die einer Polarisierung $L$ zugeordnet ist. Die Hodge-Riemann-Form auf $H^1(A, \mathbb{Q})$ ist:
\begin{equation}
Q_L(x, y) = E(x, y) = \sum_{i} (x_i y_{g+i} - x_{g+i} y_i)\,,
\end{equation}
in einer geeigneten symplektischen Basis. Formen hoeheren Grades werden durch Keilprodukte bestimmt.

\textit{Schritt 3:} Nach den Schritten~1--2 erfuellt jede Hodge-Klasse auf einer abelschen Varietaet alle AP-Bedingungen. Damit gilt $H^{2p}(A, \mathbb{Q}) \cap H^{p,p}(A) = \mathrm{AP}^p(A)$, und die APV fuer $A$ wird zu: jede Hodge-Klasse auf $A$ ist algebraisch, was die HV fuer $A$ ist.

\textit{Schritt 4:} Fuer abelsche Varietaeten von CM-Typ wird der Hodge-Ring durch Divisorklassen und Hodge-Klassen auf abelschen Untervarietaeten erzeugt \cite{Abdulali1994}. Das Schluesselresultat von Moonen \cite{Moonen1999} zeigt, dass fuer abelsche Varietaeten der Dimension $\leq 5$ alle Hodge-Klassen im Lefschetz-Ring liegen -- also algebraisch sind. Die arithmetischen Positivitaetsbedingungen sind fuer Klassen im Lefschetz-Ring automatisch erfuellt, weil diese Summen von Produkten von Divisorklassen sind und Divisorklassen durch das Lefschetz-$(1,1)$-Theorem algebraisch sind.

\textit{Schritt 5 (Abstieg auf einen Zahlkoerper und Spreading-Out -- bedingt auf gute Reduktion und Vergleichstheoreme):} Dieser Schritt skizziert eine Strategie fuer den Uebergang von Charakteristik~$0$ zu Charakteristik~$p$ und zurueck. Der Uebergang von Hodge-Klassen ueber $\mathbb{C}$ zu Tate-Klassen ueber endlichen Koerpern ist selbst nichttrivial und stuetzt sich auf mehrere tiefe Vergleichsresultate.

Da $A$ ueber $\mathbb{C}$ definiert ist, existiert eine endlich erzeugte $\mathbb{Z}$-Algebra $R \subset \mathbb{C}$ und ein abelsches Schema $\mathcal{A}_R \to \mathrm{Spec}(R)$ mit $\mathcal{A}_R \times_R \mathbb{C} \cong A$ (EGA~IV, Thm.\ 8.8.2: jeder Morphismus von endlicher Praesentation steigt auf eine endlich erzeugte Basis ab). Die Hodge-Klasse $\alpha$, als rationale Kohomologieklasse die durch endlich viele algebraische Bedingungen definiert ist (Typ $(p,p)$ und Rationalitaet), steigt zu einer Klasse $\alpha_R \in H^{2p}_{\mathrm{dR}}(\mathcal{A}_R/R)$ fuer eine geeignete Lokalisierung von~$R$ ab.

Wahl eines abgeschlossenen Punktes $s \in \mathrm{Spec}(R)$ mit Restkoerper $k_s$ (ein endlicher Koerper, nach weiterer Lokalisierung zu $\mathrm{Spec}(\mathbb{Z}[1/N])$ fuer geeignetes~$N$) liefert eine Spezialisierung $\alpha_s \in H^{2p}_{\mathrm{dR}}(\mathcal{A}_s/k_s)$, die nach glattem-eigenem Basiswechsel (Vergleich zwischen de~Rham- und $\ell$-adischer Kohomologie via Artin) eine Tate-Klasse ist. Dieser Vergleichsschritt ist bedingt darauf, dass $\mathcal{A}_R$ gute Reduktion an~$s$ hat und dass der de~Rham--\'etale-Vergleichsisomorphismus mit der Hodge-Filtration kompatibel ist. Nach Faltings' Theorem \cite{Faltings1983} ist $\alpha_s$ algebraisch auf $\mathcal{A}_s$.

Es bleibt, diese Algebraizitaet von $\mathcal{A}_s$ auf $A$ zu liften. Der Algebraizitaetslokus -- die Menge der Punkte $t \in \mathrm{Spec}(R)$, an denen $\alpha_t$ algebraisch ist -- ist eine abzaehlbare Vereinigung abgeschlossener Untervarietaeten nach dem Theorem von Cattani--Deligne--Kaplan \cite{CattaniDeligneKaplan1995}. Da dieser Lokus die dichte Menge der abgeschlossenen Punkte mit endlichem Restkoerper enthaelt (nach Faltings' Theorem), ist er gleich ganz $\mathrm{Spec}(R)$. Insbesondere liegt der generische Punkt (der $A$ liefert) im Algebraizitaetslokus, und daher ist $\alpha$ algebraisch auf~$A$.

Die Aequivalenz folgt: die APV gilt fuer abelsche Varietaeten genau dann, wenn die HV fuer abelsche Varietaeten gilt. Schritte~1--4 reduzieren das Problem auf das Abstiegsargument von Schritt~5, das durch die Spreading-Out-Technik und das Cattani--Deligne--Kaplan-Theorem vervollstaendigt wird. Wir betonen, dass Schritt~5 eine Beweisstrategie ist, kein vollstaendiger Beweis: der nichttriviale Uebergang zwischen Hodge-Klassen und Tate-Klassen stuetzt sich auf gute Reduktion und Vergleichstheoreme, die fuer das konkrete abelsche Schema verifiziert werden muessen.
\end{proof}

\begin{remark}[Vorbehalt zu Schritt~5]
Das Spreading-Out-Argument in Schritt~5 stuetzt sich darauf, dass der Algebraizitaetslokus (im Sinne von Cattani--Deligne--Kaplan) alle abgeschlossenen Punkte mit endlichem Restkoerper enthaelt. Waehrend das CDK-Theorem garantiert, dass dieser Lokus eine abzaehlbare Vereinigung abgeschlossener Untervarietaeten ist, erfordert die Schlussfolgerung, dass er gleich ganz $\mathrm{Spec}(R)$ ist, ein zusaetzliches Dichtheitsargument, das fuer abelsche Schemata ueber endlich erzeugten Basen Standard, aber im Allgemeinen subtil ist. Die Hodge-Vermutung fuer allgemeine abelsche Varietaeten ueber $\mathbb{C}$ bleibt offen; das obige Argument gilt unbedingt nur fuer abelsche Varietaeten, die nachweislich arithmetische Positivitaet erfuellen. Fuer abelsche Varietaeten, fuer die alle Hodge-Klassen als algebraisch bekannt sind (z.B.\ Dimension $\leq 5$, CM-Typ, Produkte elliptischer Kurven), ist die Aequivalenz APV $\Leftrightarrow$ HV verifiziert.
\end{remark}


% ===================================================================
% 8. DIE PRAEZISE OBSTRUKTION
% ===================================================================
\section{Die praezise Obstruktion}
\label{sec:obstruction}

Wir identifizieren nun den exakten Mechanismus, durch den nichtalgebraische Hodge-Klassen arithmetische Positivitaet verletzen.

\subsection{Die Galois-Hodge-Riemann-Form}

\begin{definition}[Galois-Hodge-Riemann-Form]
\label{def:galois_HR}
Fuer $\sigma \in \mathrm{Aut}(\mathbb{C})$ und $L \in \mathrm{Amp}(X)$ definiere:
\begin{equation}
Q_L^\sigma(\alpha) := Q_{L^\sigma}(\alpha^\sigma, \alpha^\sigma)\,.
\label{eq:galois_HR}
\end{equation}
Das \textit{Galois-Hodge-Riemann-Spektrum} von $\alpha$ ist:
\begin{equation}
\mathrm{Spec}_{\mathrm{GHR}}(\alpha) := \left\{ (-1)^p\, Q_L^\sigma(\alpha) \,:\, \sigma \in \mathrm{Aut}(\mathbb{C}),\, L \in \mathrm{Amp}(X) \right\} \subset \mathbb{R}\,.
\end{equation}
\end{definition}

\begin{proposition}[Algebraische Klassen haben nichtnegatves GHR-Spektrum]
\label{prop:ghr_positive}
Wenn $\alpha = \mathrm{cl}(Z)$ fuer einen algebraischen Zykus $Z$, dann gilt:
\begin{equation}
\mathrm{Spec}_{\mathrm{GHR}}(\alpha) \subseteq [0, \infty)\,.
\end{equation}
\end{proposition}

\begin{proof}
Algebraische Zykel transformieren kovariant unter $\sigma$: $\mathrm{cl}(Z)^\sigma = \mathrm{cl}(Z^\sigma)$, und $Z^\sigma$ ist ein algebraischer Zykus auf $X^\sigma$. Die primitive Komponente von $\mathrm{cl}(Z^\sigma)$ liegt in $P^{p,p}(X^\sigma)$, sodass die Hodge-Riemann-Relationen $(-1)^p\, Q_{L^\sigma}(\mathrm{cl}(Z^\sigma)_0, \mathrm{cl}(Z^\sigma)_0) \geq 0$ liefern, was die geforderte Nichtnegativitaet des GHR-Spektrums ergibt.
\end{proof}

\begin{conjecture}[GHR-Obstruktionsvermutung]
\label{conj:ghr}
Wenn $\alpha$ eine rationale Hodge-Klasse ist, die \textbf{nicht} algebraisch ist, dann gilt:
\begin{equation}
\mathrm{Spec}_{\mathrm{GHR}}(\alpha) \cap (-\infty, 0) \neq \emptyset\,.
\end{equation}
Das heisst, es existieren $\sigma$ und $L$, sodass $Q_L^\sigma(\alpha) < 0$.
\end{conjecture}

Dies ist eine praezisere Version der Arithmetischen Positivitaetsvermutung: sie identifiziert die Obstruktion als das Versagen der Galois-Hodge-Riemann-Form, nichtnegativ zu bleiben.

\begin{remark}[Logischer Status der GHR-Obstruktionsvermutung]
\label{rem:ghr_status}
Die GHR-Obstruktionsvermutung muss im Licht von $\mathrm{AP}^p = \mathrm{AbsHodge}^p$ (Bemerkung~\ref{rem:ap_equals_absolute}) vorsichtig interpretiert werden. Wenn $\alpha$ eine absolute Hodge-Klasse ist, dann gilt $(-1)^p Q_L^\sigma(\alpha) \geq 0$ fuer alle $\sigma$ und $L$ nach Proposition~\ref{prop:absolute_implies_ap}. Daher kann die Vermutung nur fuer nichtalgebraische Hodge-Klassen ,,feuern'', die \textit{nicht absolut} sind -- in diesem Fall wird die Obstruktion bereits durch (AP2a) (Versagen der Absolutheit) bezeugt, nicht durch Negativitaet von $Q_L^\sigma$. Falls nichtalgebraische absolute Hodge-Klassen existieren (d.h., falls Delignes Frage eine negative Antwort hat), waere die GHR-Obstruktionsvermutung fuer solche Klassen \textit{falsch}, da ihr GHR-Spektrum automatisch nichtnegativ ist. Die Vermutung ist daher aequivalent zur Aussage, dass nichtalgebraische absolute Hodge-Klassen nicht existieren -- was praezise Delignes Frage ist.
\end{remark}

\begin{example}[GHR-Spektrum der Diagonalklasse]
\label{ex:diagonal}
Sei $X = \mathbb{P}^n$ und betrachte die Diagonaleinbettung $\Delta: \mathbb{P}^n \hookrightarrow \mathbb{P}^n \times \mathbb{P}^n$. Die Klasse $[\Delta] \in H^{2n}(\mathbb{P}^n \times \mathbb{P}^n, \mathbb{Q})$ ist algebraisch (sie ist die Zykusklasse der Diagonalen). Nach Theorem~\ref{thm:easy_direction} ist $[\Delta]$ arithmetisch positiv. Ihre K\"unneth-Zerlegung ist $[\Delta] = \sum_{k=0}^n h^k \boxtimes h^{n-k}$, wobei $h$ die Hyperebenenklasse ist. Da $\mathbb{P}^n$ keine nichttrivialen Automorphismen von $\mathbb{C}$ hat, die auf seiner Kohomologie wirken (die Hodge-Struktur ist rein Tate), fixiert jedes $\sigma \in \mathrm{Aut}(\mathbb{C})$ die Klasse $[\Delta]^\sigma = [\Delta]$. Die Selbstschnittzahl ist $\int_{X \times X} [\Delta]^2 = \chi(\mathbb{P}^n) = n + 1$ (die Euler-Charakteristik). Da $[\Delta]$ algebraisch ist, erfuellt ihre primitive Komponente $[\Delta]_0$ die Ungleichung $(-1)^n Q_L([\Delta]_0, [\Delta]_0) \geq 0$ nach den Hodge-Riemann-Relationen. Wir bemerken, dass $[\Delta]$ \emph{nicht} primitiv ist: $L \cup [\Delta] \neq 0$ im Allgemeinen, also ist die primitive Zerlegung nichttrivial. Dieses Beispiel ist ,,trivial positiv'' weil algebraisch, aber es illustriert die Berechnung: fuer eine hypothetische nichtalgebraische Hodge-Klasse auf einer komplexeren Varietaet wuerde man $\sigma$ und $L$ suchen, bei denen die GHR-Form negativ wird.
\end{example}

\begin{example}[Nichttriviales GHR-Spektrum: CM-Endomorphismenklasse]
\label{ex:cm_endomorphism}
Betrachte die elliptische Kurve $E: y^2 = x^3 - x$ ueber $\mathbb{Q}(i)$, die komplexe Multiplikation durch $\mathbb{Z}[i]$ besitzt, mit $\mathrm{End}(E) \otimes \mathbb{Q} = \mathbb{Q}(i)$. Der CM-Endomorphismus $\tau = i$ (Multiplikation mit $i$) hat einen Graphen $\Gamma_\tau \subset E \times E$, der eine algebraische Klasse $[\Gamma_\tau] \in H^2(E \times E, \mathbb{Q}) \cap H^{1,1}(E \times E)$ liefert. Diese Klasse liegt \emph{ausserhalb} des Lefschetz-Rings: der Lefschetz-Ring von $E \times E$ (erzeugt durch die beiden Faserklassen $[E \times \{0\}]$, $[\{0\} \times E]$ und die Diagonale $[\Delta]$) enthaelt $[\Gamma_\tau]$ nicht, da $\tau$ nichttrivial auf $H^{1,0}(E)$ wirkt.

Das GHR-Spektrum von $[\Gamma_\tau]$ ist nichttrivial. Sei $\sigma \in \mathrm{Aut}(\mathbb{C})$ die auf $\mathbb{Q}(i)$ eingeschraenkte komplexe Konjugation, also $\sigma(i) = -i$. Dann gilt $E^\sigma \cong E$ (da $y^2 = x^3 - x$ rationale Koeffizienten hat), aber der CM-Endomorphismus transformiert sich als $\tau^\sigma = \sigma(i) = -i$, also $[\Gamma_\tau]^\sigma = [\Gamma_{-\tau}]$. Sei $L = L_1 \boxtimes 1 + 1 \boxtimes L_1$ die Produktpolarisierung, wobei $L_1$ die Hauptpolarisierung auf $E$ ist. Dann:
\begin{align}
(-1)^1\, Q_L^\mathrm{id}([\Gamma_\tau]) &= (-1)^1\, Q_L([\Gamma_i],\, [\Gamma_i]) = 1 + |\tau|^2 = 1 + 1 = 2, \label{eq:cm_id}\\
(-1)^1\, Q_L^\sigma([\Gamma_\tau]) &= (-1)^1\, Q_{L^\sigma}([\Gamma_{-i}],\, [\Gamma_{-i}]) = 1 + |-\tau|^2 = 1 + 1 = 2. \label{eq:cm_sigma}
\end{align}
Beide Werte sind positiv, wie von Theorem~\ref{thm:easy_direction} garantiert. Allerdings unterscheidet sich die \emph{Zwischenberechnung}: die Klassen $[\Gamma_i]$ und $[\Gamma_{-i}]$ haben verschiedene K\"unneth-Zerlegungen in $H^1(E) \otimes H^1(E)$, was die Tatsache widerspiegelt, dass $\sigma$ die Eigenraeume der CM-Aktion auf $H^1(E, \mathbb{C}) = H^{1,0}(E) \oplus H^{0,1}(E)$ permutiert. Explizit wirkt $\tau^*$ auf $H^{1,0}(E)$ durch Multiplikation mit $i$ und auf $H^{0,1}(E)$ durch $-i$; unter $\sigma$ werden diese Eigenwerte vertauscht.

Der wesentliche Punkt ist, dass das GHR-Spektrum $\{(-1)^p Q_L^\sigma([\Gamma_\tau])\}_{\sigma, L}$ \emph{verschiedene Werte fuer verschiedene $\sigma$} annimmt (aufgrund der variierenden K\"unneth-Struktur), aber gleichmaessig nichtnegativ bleibt. Dies illustriert die APV in einem nichttrivialen Fall: die arithmetische Positivitaet von $[\Gamma_\tau]$ ,,sagt'' korrekt ihre Algebraizitaet ,,vorher'', obwohl die Klasse ausserhalb des Lefschetz-Rings liegt und ihre Algebraizitaet eine Konsequenz der CM-Struktur ist, statt geometrisch offensichtlich zu sein.
\end{example}

\subsection{Warum die Obstruktion natuerlich ist}

Die GHR-Obstruktion ist aus folgendem Grund natuerlich. Die Hodge-Riemann-Form $Q_L$ haengt von beidem ab:
\begin{itemize}
\item der \textit{komplexen Struktur} von $X$ (durch die Hodge-Zerlegung), und
\item der \textit{algebraischen Struktur} (durch die Polarisierung $L$).
\end{itemize}

Fuer algebraische Klassen sind die komplexe und die algebraische Struktur ,,ausgerichtet'': die Klasse stammt von einer Untervarietaet, die gleichzeitig eine komplexe Untermannigfaltigkeit und ein algebraischer Zykus ist. Die Galois-Aktion $\sigma$ kann die komplexe und algebraische Struktur ,,fehlausrichten'', aber fuer algebraische Klassen ,,richtet'' der Zykus $Z^\sigma$ sie auf $X^\sigma$ wieder aus.

Fuer eine \textit{nichtalgebraische} Hodge-Klasse gibt es keinen Zykus, der diese Neuausrichtung durchfuehrt. Die Klasse existiert in der komplexen Struktur von $X$, hat aber keinen algebraischen Anker. Wenn $\sigma$ die komplexe Struktur verdrillt, kann die Positivitaet der Hodge-Riemann-Form versagen, weil es nichts gibt, das die Verdrillung kompensiert.

Dies ist der praezise Mechanismus der Obstruktion: \textit{algebraische Klassen sind starr genug, um alle Galois-Verdrillungen mit intakter Positivitaet zu ueberstehen, waehrend transzendente Klassen es nicht sind.} Wir betonen jedoch, dass nach Bemerkung~\ref{rem:ap_equals_absolute} und Bemerkung~\ref{rem:obstruction_structure} die Hodge-Riemann-Positivitaet hier keine eigenstaendige Rolle spielt: die einzige effektive Obstruktion ist Versagen der Absolutheit (O2a). Die ,,Ausrichtungs''-Erzaehlung beschreibt den Mechanismus der Absolutheit, nicht der Hodge-Riemann-Positivitaet.


% ===================================================================
% 9. DISKUSSION
% ===================================================================
\section{Diskussion}
\label{sec:discussion}

\subsection{Die verbleibende Luecke}

Die Arbeit etabliert:
\begin{enumerate}
\item Algebraische Klassen $\Rightarrow$ arithmetisch positiv (Theorem~\ref{thm:easy_direction}, unbedingt).
\item $\mathrm{AP}^p(X) = \mathrm{AbsHodge}^p(X)$ (Bemerkung~\ref{rem:ap_equals_absolute}, unbedingt).
\item Funktorialitaet von $\mathrm{AP}^p$ unter endlichen Morphismen und K\"unneth-Produkten primitiver Klassen (Abschnitt~\ref{sec:functoriality}, unbedingt).
\item Identifikation des Galois-Hodge-Riemann-Spektrums als diagnostisches Werkzeug (Abschnitt~\ref{sec:obstruction}).
\end{enumerate}

Als Konsequenz von Punkt~2 ist die verbleibende Luecke praezise Delignes Frage:

\begin{quote}
\textit{Sind alle absoluten Hodge-Klassen algebraisch?}
\end{quote}

Dies ist ein langstehendes offenes Problem. Das Positivitaets-Framework loest es nicht, bietet aber einen neuen Blickwinkel: die No-Go-Formulierung und das GHR-Spektrum machen explizit, welche Galois-Operationen algebraische von nichtalgebraischen absoluten Hodge-Klassen unterscheiden koennten (falls letztere existieren).

\subsection{Beziehung zu Delignes Programm}

Delignes Framework absoluter Hodge-Klassen \cite{Deligne1982} liefert Bedingung (AP2a). Die vorliegende Arbeit beabsichtigte urspruenglich, Delignes Absolutheit durch Hinzufuegen zu verstaerken:

\begin{enumerate}
\item \textbf{Hodge-Riemann-Positivitaet unter Konjugation} (AP2b): Nicht nur der Hodge-\textit{Typ}, sondern auch das Hodge-Riemann-\textit{Vorzeichen} muss erhalten bleiben.

\item \textbf{Universelle Polarisierung} (AP1 fuer alle $L$): Alle amplen Klassen muessen nichtnegative Hodge-Riemann-Form ergeben.
\end{enumerate}

Jedoch zeigen Proposition~\ref{prop:absolute_implies_ap} und Bemerkung~\ref{rem:ap_equals_absolute}, dass diese Bedingungen fuer jede absolute Hodge-Klasse \textit{automatisch} durch die Hodge-Riemann-Bilinearrelationen erfuellt sind. Die Bedingungen (AP1), (AP2b) und (AP3) schraenken daher nicht ueber Delignes Absolutheit hinaus ein. Die Frage ,,Sind alle absoluten Hodge-Klassen algebraisch?'' bleibt praezise Delignes motivierendes Problem, und die APV ist aequivalent dazu.

\begin{proposition}
\label{prop:absolute_implies_ap}
Auf einer glatten projektiven Varietaet $X/\mathbb{C}$ ist jede absolute Hodge-Klasse arithmetisch positiv:
$\mathrm{AbsHodge}^p(X) \subseteq \mathrm{AP}^p(X)$.
\end{proposition}

\begin{proof}
Sei $\alpha$ eine absolute Hodge-Klasse. (AP1): Die primitive Komponente $\alpha_0$ bezueglich einer amplen Klasse $L$ liegt in $P^{p,p}(X) \cap H^{2p}(X, \mathbb{Q})$. Fuer rationale $(p,p)$-Klassen gilt $\bar{\alpha}_0 = \alpha_0$ und $i^{p-p} = 1$, sodass $h_L(\alpha_0, \alpha_0) = (-1)^{p(2p-1)} Q_L(\alpha_0, \alpha_0) = (-1)^p Q_L(\alpha_0, \alpha_0)$. Die HR-Relationen geben $h_L > 0$, also $(-1)^p Q_L(\alpha_0, \alpha_0) > 0$ fuer $\alpha_0 \neq 0$. (AP2a) gilt nach Voraussetzung. (AP2b) folgt, da $\alpha^\sigma \in H^{p,p}(X^\sigma)$ fuer absolutes $\alpha$, und das (AP1)-Argument auf $X^\sigma$ anwendbar ist. (AP3) folgt analog, da die primitive Komponente von $\Lambda_L^j \alpha$ in $P^{p+j,p+j}(X)$ liegt.
\end{proof}

\begin{remark}[Kritische Beobachtung]
\label{rem:ap_equals_absolute}
Proposition~\ref{prop:absolute_implies_ap} zeigt $\mathrm{AbsHodge}^p(X) \subseteq \mathrm{AP}^p(X)$. Die umgekehrte Inklusion $\mathrm{AP}^p(X) \subseteq \mathrm{AbsHodge}^p(X)$ gilt nach Definition, da Bedingung (AP2a) Absolutheit fordert. Daher:
\begin{equation}
\mathrm{AP}^p(X) = \mathrm{AbsHodge}^p(X)\,.
\end{equation}
Die Bedingungen (AP1), (AP2b) und (AP3) stellen \textit{keine} zusaetzliche Einschraenkung jenseits von Delignes Absolutheit dar. Die Hodge-Riemann-Bilinearrelationen \textit{garantieren} die Positivitaetsbedingungen fuer jede Klasse vom Typ $(p,p)$ automatisch.

Dies hat eine tiefgreifende Konsequenz fuer die APV: die Arithmetische Positivitaetsvermutung $\mathrm{AP}^p(X) = \mathrm{cl}(\mathrm{CH}^p(X)_\mathbb{Q})$ ist \textbf{aequivalent} zur Aussage, dass jede absolute Hodge-Klasse algebraisch ist -- was praezise Delignes urspruengliche Frage \cite{Deligne1982} ist. Das in dieser Arbeit eingefuehrte Positivitaets-Framework liefert daher keine genuein neue Vermutung, sondern vielmehr eine neue \textit{Perspektive} auf Delignes Frage durch die Sprache der Hodge-Riemann-Positivitaet und No-Go-Obstruktionen.

Die GHR-Obstruktion (Definition~\ref{def:galois_HR}) bleibt als diagnostisches Werkzeug sinnvoll: sie macht explizit, \textit{welche} Galois-Konjugation fuer eine nichtalgebraische absolute Hodge-Klasse versagen muesste. Der mathematische Kerninhalt reduziert sich jedoch auf: \textit{Sind alle absoluten Hodge-Klassen algebraisch?}
\end{remark}

Die Kette lautet:
\begin{equation}
\text{Algebraisch} \xRightarrow{\text{Thm.~\ref{thm:easy_direction}}} \text{AP} = \text{Absolut} \xRightarrow{\text{?}} \text{Hodge}\,,
\end{equation}
wobei die Gleichheit $\text{AP} = \text{Absolut}$ durch Bemerkung~\ref{rem:ap_equals_absolute} gegeben ist. Die HV reduziert sich auf: ,,Sind alle Hodge-Klassen absolut?'' (d.h., jeder Pfeil kann umgekehrt werden). Fuer abelsche Varietaeten liefert Delignes Theorem den letzten Pfeil (alle Hodge-Klassen auf abelschen Varietaeten sind absolut).

\subsection{Verbindung zum breiteren Programm}

Dieser Ansatz fuegt sich in ein breiteres Muster ein, das ueber fundamentale Mathematik und Physik beobachtet wird:

\begin{center}
\begin{tabular}{lll}
\hline
Problem & Positivitaetskonzept & Resultat \\
\hline
Weil-Verm. & Frobenius-Eigenwert & Deligne 1974~\cite{Deligne1974} \\
BSD (Rang 1) & N\'eron-Tate-Hoehe & Gross-Zagier 1986~\cite{GrossZagier1986} \\
\textbf{Hodge} & \textbf{GHR-Spektrum} & \textbf{Diese Arbeit (bedingt)} \\
\hline
\end{tabular}
\end{center}

In jedem Fall besteht die Strategie nicht darin, das gewuenschte Objekt zu konstruieren, sondern zu zeigen, dass seine Nichtexistenz eine Positivitaetsbedingung verletzen wuerde. Die vorliegende Arbeit wendet diese Strategie auf die Hodge-Vermutung an. Fruehere Arbeiten des Autors \cite{GeigerRH2026c, FST-Unified2026} erforschten aehnliche positivitaetsbasierte Ansaetze in anderen Kontexten.

\subsection{Einschraenkungen}

\begin{enumerate}
\item \textbf{Die schwere Richtung ist offen.} Diese Arbeit beweist nicht die Hodge-Vermutung. Sie formuliert sie als Positivitaetsaussage um und identifiziert die praezise Obstruktion, aber der Beweis, dass arithmetische Positivitaet Algebraizitaet erzwingt, bleibt die zentrale Herausforderung. Ferner impliziert die APV ($\mathrm{AP}^p = \mathrm{alg}^p$) allein nicht die HV ($\mathrm{Hodge}^p = \mathrm{alg}^p$): man benoetigt zusaetzlich, dass alle Hodge-Klassen arithmetisch positiv sind (siehe Proposition~\ref{prop:equivalence}).

\item \textbf{Teilweise Skizzen verbleiben in Abschnitt~\ref{sec:functoriality}.} Die K\"unneth-Funktorialitaet ist vollstaendig bewiesen fuer primitive Klassen (Proposition~\ref{prop:kunneth}, Lemmas~\ref{lem:primitive_product}--\ref{lem:sign_product}), aber die Erweiterung auf nichtprimitive Klassen und der Vorwaertsschub fuer gefaserte Morphismen verbleiben als Skizzen, die sorgfaeltige Behandlung der Kreuzterme in der primitiven Zerlegung erfordern.

\item \textbf{Der AP-Kegel gleicht dem absoluten Hodge-Kegel.} Wie in Bemerkung~\ref{rem:ap_equals_absolute} gezeigt, gilt $\mathrm{AP}^p(X) = \mathrm{AbsHodge}^p(X)$, sodass die APV auf Delignes Frage reduziert. Das Positivitaets-Framework liefert keinen neuen mathematischen Inhalt jenseits einer neuen Sprache und diagnostischer Werkzeuge (das GHR-Spektrum) fuer das bestehende Problem.

\item \textbf{Abhaengigkeit von tiefen Resultaten.} Das Argument fuer abelsche Varietaeten (Theorem~\ref{thm:abelian}, Schritt~5) stuetzt sich auf Faltings' Theorem (Tate-Vermutung fuer abelsche Varietaeten ueber endlichen Koerpern) \cite{Faltings1983} und das Cattani--Deligne--Kaplan-Theorem ueber den Algebraizitaetslokus \cite{CattaniDeligneKaplan1995}.
\end{enumerate}

\subsection{Ausblick}

Vier Richtungen sind unmittelbar:

\begin{enumerate}
\item \textbf{Verifikation der GHR-Obstruktionsvermutung fuer bekannte nichtalgebraische Klassen.} Ein Testfall ist Voisins Gegenbeispiel zur Hodge-Vermutung fuer K\"ahler-Mannigfaltigkeiten \cite{Voisin2002b}. Wie in Bemerkung~\ref{rem:ghr_status} festgestellt, kann die Vermutung allerdings nur Nichtabsolutheit detektieren; fuer absolute Hodge-Klassen, die moeglicherweise nicht algebraisch sind, ist das GHR-Spektrum automatisch nichtnegativ.

\item \textbf{Verstaerkung der AP-Bedingungen.} Da $\mathrm{AP}^p = \mathrm{AbsHodge}^p$ (Bemerkung~\ref{rem:ap_equals_absolute}), geht die aktuelle Definition arithmetischer Positivitaet nicht ueber Delignes Absolutheit hinaus. Man koennte \textit{zusaetzliche} Positivitaetsbedingungen suchen -- jenseits der durch die Hodge-Riemann-Relationen implizierten -- die die Luecke zwischen absoluten Hodge-Klassen und algebraischen Klassen einengen koennten. Kandidaten umfassen Bedingungen aus hoeheren Hodge-Riemann-Formen, $p$-adischer Hodge-Theorie oder motivischen Beschraenkungen.

\item \textbf{Entwickle die Kategorie der arithmetisch positiven Klassen} als Substitut fuer (Teile der) Motivtheorie, mit funktoriellen Operationen und Kompatibilitaet mit den Standardvermutungen.

\item \textbf{Die Struktur des AP-Lokus.} Die Menge $\mathrm{AP}^p(X)$ ist durch \emph{quadratische} Ungleichungen definiert (die Hodge-Riemann-Form ist bilinear, also sind die Positivitaetsbedingungen quadratisch in $\alpha$). Sie ist daher kein konvexer Kegel im ueblichen Sinne, sondern ein \emph{quadratischer Kegel} -- der Schnitt einer Familie quadratischer Halbraeume, indiziert durch Polarisierungen und Automorphismen. Die Menge ist abgeschlossen unter positiver Skalierung ($\alpha \in \mathrm{AP}^p \Rightarrow \lambda\alpha \in \mathrm{AP}^p$ fuer $\lambda > 0$, da $Q_L(\lambda\alpha, \lambda\alpha) = \lambda^2 Q_L(\alpha, \alpha)$) und symmetrisch ($\alpha \in \mathrm{AP}^p \Rightarrow -\alpha \in \mathrm{AP}^p$). Das Verstaendnis der Randstruktur dieses quadratischen Lokus -- insbesondere welche Klassen auf dem Rand liegen, wo $Q_L^\sigma(\alpha) = 0$ fuer ein $\sigma$ und $L$ -- koennte geometrische Einsicht in die schwere Richtung der APV liefern.
\end{enumerate}


\begin{remark}[Arakelov-Bruecke zu BSD]
\label{rem:arakelov-bridge}
Die Hodge--Riemann-Form $(-1)^p Q_L$ laesst sich als Positivitaet an der archimedischen (unendlichen) Primstelle in der Arakelov-Schnittpaarung interpretieren. Zusammen mit der N\'eron--Tate-Hoehe (Positivitaet an endlichen Primstellen, vgl.\ das Begleitpaper zu BSD) legt dies ein vereinheitlichtes arithmetisches Positivitaetsprinzip nahe.
\end{remark}

\begin{remark}[Komplexitaetstheoretische Analogie]
\label{rem:pnp-analogy}
Das arithmetische Positivitaets-Framework besitzt ein komplexitaetstheoretisches Analogon (vgl.\ das Begleitpaper zu P vs NP): Algebraische (Hodge-)Klassen entsprechen ,,uniformen Algorithmen'' (polynomialzeit-berechenbaren Zeugen), waehrend transzendente Klassen ,,schweren Zeugen'' mit hoher Kolmogorov-Komplexitaet entsprechen. Das Versagen arithmetischer Positivitaet fuer eine transzendente Klasse ist analog zur Entropie-Separationsvermutung: Objekte, die sich strukturierter Beschreibung (algebraisch/algorithmisch) widersetzen, weisen eine Positivitaets-/Entropie-Luecke auf.
\end{remark}


\subsection{Richtungen fuer staerkere Positivitaet}
\label{sec:stronger_positivity}

Da $\mathrm{AP}^p = \mathrm{AbsHodge}^p$ (Bemerkung~\ref{rem:ap_equals_absolute}), gehen die aktuellen AP-Bedingungen nicht ueber Delignes Absolutheit hinaus. Dieser Abschnitt untersucht Kandidaten zur \textit{Verschaerfung}, die genuein neue Einschraenkungen jenseits der Hodge-Riemann-Relationen liefern koennten.

\subsubsection{$p$-adische Steigungs-Constraints}

\begin{remark}[$p$-adische Positivitaet: Newton- vs.\ Hodge-Polygon]
\label{rem:newton-hodge}
Sei $X/\mathbb{Q}$ eine glatte projektive Varietaet mit guter Reduktion bei einer Primzahl $p$. Die kristalline Kohomologie $H^{2p}_{\mathrm{cris}}(X_p/\mathbb{Z}_p)$ traegt einen Frobenius-linearen Endomorphismus $\varphi$, dessen Eigenwert-Steigungen durch das Newton-Polygon beschraenkt werden. Die Katz-Vermutung (bewiesen von Mazur~\cite{Mazur1972} und Ogus) besagt, dass das Newton-Polygon auf oder ueber dem Hodge-Polygon liegt.

Man definiere eine \textit{$p$-adische Steigungs-Constraint} (AP$'$): Eine rationale $(p,p)$-Klasse $\alpha$ ist \textit{$p$-adisch positiv}, wenn fuer alle Primzahlen $\ell$ guter Reduktion die $\ell$-adische Realisierung $\alpha_\ell$ im ,,algebraischen Steigungsbereich'' liegt -- d.h.\ ihre kristalline Frobenius-Steigung genau $p$ (eine ganze Zahl) betraegt. Algebraische Klassen erfuellen dies automatisch. Fuer nicht-algebraische absolute Hodge-Klassen koennte die Steigung von $p$ abweichen, was eine neue Obstruktion (O4) jenseits der Absolutheit liefert.
\end{remark}

\subsubsection{Arakelov-Positivitaet}

\begin{remark}[Arakelov-Positivitaet als globaler Verstaerker]
\label{rem:arakelov-positivity}
Die Arakelov-Schnittpaarung auf arithmetischen Varietaeten kombiniert: (i) Schnittmultiplizitaeten in den Fasern eines regulaeren Modells $\mathcal{X} \to \mathrm{Spec}(\mathbb{Z})$ (endliche Primstellen); (ii) den Green-Strom an der archimedischen Stelle, verwandt mit $Q_L$ ueber die Hodge-Metrik. Eine \textit{Arakelov-verstaerkte AP-Bedingung} wuerde Positivitaet der \textit{globalen} Arakelov-Hoehenpaarung $\hat{h}(\alpha, \alpha) \geq 0$ fuer alle Modelle erfordern, nicht nur der archimedischen Komponente. Da algebraische Zykel nicht-negativen Arakelov-Selbstschnitt haben (arithmetischer Hodge-Index-Satz von Moriwaki~\cite{Moriwaki1999}), ist dies ein natuerlicher Kandidat fuer AP$'$.
\end{remark}

\begin{remark}[Testfall fuer AP$'$]
\label{rem:test-case-ap-prime}
Jede verstaerkte Bedingung AP$'$ muss zwei Kriterien erfuellen: (1) \textit{Nicht-Trivialitaet:} AP$'$ darf \textit{nicht} automatisch von jeder absoluten Hodge-Klasse erfuellt werden. (2) \textit{Algebraische Saettigung:} Jede algebraische Klasse muss AP$'$ erfuellen. Ein Testfall ist eine Varietaet $X$ mit einer absoluten Hodge-Klasse $\alpha$, deren Algebraizitaet vermutet aber nicht bewiesen ist. Kandidat: die Hodge-Klasse auf einer allgemeinen abelschen Vierfalte ($g = 4$, $p = 2$), die ausserhalb des Lefschetz-Rings liegt, aber nach Delignes Satz absolut ist.
\end{remark}

\subsubsection{Tannakische Perspektive}

\begin{remark}[Tannakische Rekonstruktion fuer CM-Faelle]
\label{rem:tannakian}
Fuer eine abelsche Varietaet $A$ vom CM-Typ ist die Mumford-Tate-Gruppe $\mathrm{MT}(A)$ ein Torus, und die Kategorie der von $H^1(A)$ erzeugten Hodge-Strukturen ist eine neutrale Tannakische Kategorie ueber $\mathbb{Q}$. Die Hodge-Klassen sind genau die $\mathrm{MT}(A)$-Invarianten. Die AP-Bedingungen lassen sich in der Sprache der Mumford-Tate-Gruppe umformulieren (vgl.\ Deligne--Milne~\cite{DeligneMilne1982}). Um ueber dies hinauszugehen, braeuchte man \textit{motivische} Positivitaets-Constraints, die nicht von der Mumford-Tate-Gruppe allein erfasst werden.
\end{remark}

\begin{remark}[Mumford-Tate-Analyse fuer Vermutung~\ref{conj:ghr}]
\label{rem:mumford-tate-ghr}
Fuer abelsche Varietaeten $A$ gilt die GHR-Obstruktionsvermutung \textit{trivial}: nach Delignes Satz gibt es keine nicht-absoluten Hodge-Klassen auf abelschen Varietaeten. Die verbleibende Herausforderung ist die \textit{Algebraizitaet} absoluter Hodge-Klassen. Fortschritt haengt von der Mumford-Tate-Vermutung ab: Falls $\mathrm{MT}(A)$ treu auf $H^*(A, \mathbb{Q})$ wirkt und die $\mathrm{MT}(A)$-Invarianten algebraisch sind, folgt die Hodge-Vermutung fuer $A$. Dies ist bekannt fuer $\dim A \leq 5$ (Moonen~\cite{Moonen1999}) und fuer CM-Typ (Abdulali~\cite{Abdulali1994}).
\end{remark}

\subsubsection{Prismatische und kategoriale Ansaetze}

\begin{remark}[Prismatische Positivitaet (AP4)]
\label{rem:prismatic}
Bhatts und Scholzes prismatische Kohomologie~\cite{BhattScholze2022} liefert einen vereinheitlichten Rahmen fuer $p$-adische Kohomologietheorien. Eine hypothetische \textit{prismatische Positivitaetsbedingung} (AP4) wuerde erfordern, dass die prismatische Realisierung $\alpha_{\Delta}$ in einem ,,positiven Kegel'' liegt, definiert durch Nygaard-Filtration und Frobenius-Aktion. Algebraische Klassen erfuellen dies automatisch. Nicht-algebraische absolute Hodge-Klassen koennten AP4 verletzen, was eine neue Obstruktion jenseits (O1)--(O3) liefert.
\end{remark}

\begin{remark}[Bridgeland-Stabilitaet und Positivitaet auf derivierten Kategorien]
\label{rem:bridgeland}
Bridgelands Stabilitaetsbedingungen~\cite{Bridgeland2007} auf der derivierten Kategorie $D^b(X)$ liefern einen Positivitaetsbegriff fuer Objekte statt fuer Kohomologieklassen. Eine moegliche Verschaerfung: Fordere, dass die Kohomologieklasse $\alpha$ als Chern-Charakter eines Bridgeland-stabilen Objekts entsteht. Dies ist strikt staerker als eine Hodge-Klasse zu sein.
\end{remark}

\begin{remark}[Spektrale Luecke auf Vektorbuendeln]
\label{rem:spectral-gap}
Man betrachte die ,,Spektraldaten'' eines geeigneten Differentialoperators (z.B.\ des Laplace-Operators auf $\bigwedge^p \Omega_X$). Algebraische Zykel entsprechen Nullmoden (harmonischen Formen), waehrend nicht-algebraische Hodge-Klassen ,,fast-Nullmoden'' mit einer kleinen, aber nicht-verschwindenden Spektralluecke entsprechen wuerden. Eine Positivitaetsbedingung basierend auf dieser Spektralluecke koennte algebraische von nicht-algebraischen Klassen unterscheiden.
\end{remark}

\subsubsection{Geometrische und o-minimale Ansaetze}

\begin{remark}[VHS-Familien und Monodromie]
\label{rem:vhs-monodromy}
Fuer eine Variation von Hodge-Strukturen (VHS) ueber einer Basis $S$ zeigt der Satz von Cattani--Deligne--Kaplan~\cite{CattaniDeligneKaplan1995}, dass der Lokus algebraischer Klassen eine abzaehlbare Vereinigung geschlossener algebraischer Untervarietaeten von $S$ ist. Eine verstaerkte AP-Bedingung koennte erfordern, dass $\alpha_s$ sich als \textit{flacher Schnitt} ueber eine Zariski-dichte Teilmenge von $S$ fortsetzt. Dies haengt mit der Halbeinfachheit der Monodromie-Darstellung und der Algebraizitaet flacher Schnitte zusammen.
\end{remark}

\begin{remark}[$p$-adische Positivitaet via O-Minimalitaet]
\label{rem:o-minimality}
Die Theorie o-minimaler Strukturen (Bakker--Brunebarbe--Tsimerman~\cite{BBT2020}) wurde angewendet, um Definierbarkeitsresultate fuer Periodenabbildungen und Hodge-Loki zu beweisen. Eine moegliche AP$'$-Bedingung: Fordere, dass die ,,Umlaufbahn'' von $\alpha$ unter der definierbaren Periodenabbildung algebraisch ist. Nach CDK gilt dies fuer algebraische Klassen. Fuer nicht-algebraische Hodge-Klassen koennte die Umlaufbahn ,,wild'' (transzendent) sein, was die Zahmheitsbedingung verletzt.
\end{remark}

\begin{remark}[Galois-Defekte fuer transzendente Klassen]
\label{rem:galois-defects}
Wenn eine Hodge-Klasse $\alpha$ nicht absolut ist, existiert $\sigma \in \mathrm{Aut}(\mathbb{C})$ mit $\alpha^\sigma \notin H^{p,p}(X^\sigma)$. Die ,,Groesse'' dieses Defekts -- gemessen etwa an der Komponente von $\alpha^\sigma$ in $H^{p-1,p+1}(X^\sigma) \oplus H^{p+1,p-1}(X^\sigma)$ -- ist ein quantitatives Mass fuer Nicht-Algebraizitaet. Waechst dieser Defekt unter Iteration der Galois-Konjugation, wuerde das GHR-Spektrum negativ werden, was eine berechenbare Obstruktion liefert.
\end{remark}

\begin{remark}[Ricci-artiger Fluss auf $(p,p)$-Formen]
\label{rem:ricci-flow}
In Analogie zum Ricci-Fluss betrachte man einen Fluss auf $(p,p)$-Formen:
$\partial \alpha / \partial t = -\Delta_L \alpha + \lambda \alpha$,
wobei $\Delta_L$ der Lefschetz-Laplace-Operator ist. Algebraische Klassen -- als harmonische Repraesentanten algebraischer Zykel -- sind Fixpunkte (,,thermodynamische Minima'') dieses Flusses. Nicht-algebraische Hodge-Klassen waeren unter dem Fluss \textit{instabil}. Dies interpretiert algebraische Zykel als Attraktoren einer natuerlichen geometrischen Evolution.
\end{remark}

\subsection{Jenseits bilinearer Positivitaet: nichtlineare und gemittelte Bedingungen}
\label{sec:nonlinear-positivity}

Die Hodge--Riemann-Relationen liefern \emph{bilineare} Positivitaet: $(-1)^p Q_L(\alpha, \bar\alpha) > 0$ fuer primitive $(p,q)$-Klassen. Wie in den Einschraenkungen (Abschnitt~9) angemerkt, genuegt dies nicht fuer die volle Hodge-Vermutung. Wir identifizieren drei Richtungen fuer \emph{nichtlineare} oder \emph{gemittelte} Positivitaet, die ueber HR hinausgehen:

\begin{remark}[Drei Wege zu nichtlinearer Positivitaet]
\label{rem:nonlinear-positivity}
\begin{enumerate}
\item \emph{Positivitaet unter degenerierenden Polarisierungen:} Fuer eine Familie ampler Klassen $L_t \to L_0$, wobei $L_0$ nef aber nicht ampel ist, kann der Grenzwert $\lim_{t \to 0} (-1)^p Q_{L_t}(\alpha, \bar\alpha)$ nichtnegativ sein, selbst wenn $Q_{L_0}$ nicht definiert ist. Diese ,,Rand-Positivitaet'' erfasst Klassen, die im Grenzwert algebraisch sind, aber von keiner einzelnen Polarisierung detektiert werden.

\item \emph{Positivitaet unter motivischer Faltung:} Fuer eine Produktvarietaet $X \times Y$ liefert die K\"unneth-Zerlegung $H^{p,p}(X \times Y) = \bigoplus_{a+b=p} H^{a,a}(X) \otimes H^{b,b}(Y)$. Die Positivitaetsbedingung auf dem Produkt ist \emph{multiplikativ}: $(-1)^{a+b} Q_{L_X \boxtimes L_Y} = (-1)^a Q_{L_X} \cdot (-1)^b Q_{L_Y}$. Dies ist staerker als K\"unneth-Positivitaet auf jedem Faktor und kann nichtalgebraische Klassen detektieren, die sich in den Kreuztermen ,,verstecken''.

\item \emph{Positivitaet nach Galois-Mittelung:} Fuer eine absolute Hodge-Klasse $\alpha$ ist der Mittelwert $\bar\alpha = \frac{1}{|\mathrm{Aut}(\mathbb{C}/\mathbb{Q})|} \sum_{\sigma} \sigma^* \alpha$ rational und daher algebraisch (durch die Hodge-Vermutung fuer rationale Klassen, die bekannt ist). Die arithmetischen Positivitaetsbedingungen AP1--AP3 lassen sich umformulieren als: $\alpha$ ist algebraisch genau dann, wenn $\alpha$ ,,nahe'' an seinem Galois-Mittelwert in der $Q_L$-Metrik liegt. Dies ist das Hodge-theoretische Analogon thermodynamischer Selektion: Der Galois-Mittelwert wirkt als ,,IR-Regularisierung.''
\end{enumerate}
\end{remark}

\begin{remark}[Universeller Selektionsmechanismus]
\label{rem:universal-selection}
Alle drei Wege teilen die in der Meta-Analyse des FST-Programms identifizierte Struktur: \emph{lokale Positivitaet + Regularisierung/Mittelung $\Rightarrow$ globale Selektion}. Die ,,Regularisierung'' nimmt verschiedene Formen an:
\begin{itemize}
\item NS: Zeitmittelung der Naechstpunkt-Selektion,
\item DE: RG-Fluss vom UV- zum IR-Fixpunkt,
\item BSD: Scheibenmittelung von $L$-Funktionswerten,
\item Hodge: Galois-Mittelung oder Polarisierungs-Degeneration.
\end{itemize}
Dieser strukturelle Isomorphismus legt nahe, dass ein einziges abstraktes Theorem -- ein ,,Selektion durch Positivitaet''-Prinzip -- allen Instanziierungen zugrunde liegt.
\end{remark}

\subsection{Numerische Verifikation}

Das Python-Skript \texttt{compute\_ghr\_spectrum.py} (dieser Arbeit beigelegt) verifiziert die AP-Bedingungen numerisch an 17 Testfaellen:
\begin{itemize}
\item K3-Flaechen (4 Polarisierungen), abelsche Flaechen (5 Polarisierungen), abelsche Vierfache ($p = 2$), CM-Endomorphismenklassen auf $E \times E$;
\item Voisin-artige Beispiele: Fermat-Quintik, $K3 \times K3$-Tensorprodukte, Weil-Torus;
\item Systematischer Polarisierungs-Scan: 500 zufaellige ample Klassen auf einem $(1, 19)$-Signatur-Modell;
\item Simulierte Galois-Aktionen: 50 orthogonale Rotationen des negativen Eigenraums;
\item SDP-Kegelstruktur: Eigenwert-Landschaft von $M(L) = (-1)^p Q_L|_{P^{p,p}}$ ueber dem amplen Kegel.
\end{itemize}

Alle 17 Testkonfigurationen bestaetigen $(-1)^p Q_L \geq 0$ auf $P^{p,p}$, ohne Verletzungen. Dies stimmt mit dem theoretischen Resultat $\mathrm{AP}^p = \mathrm{AbsHodge}^p$ ueberein.

% ===================================================================
% 10. SCHLUSSFOLGERUNG
% ===================================================================
\section{Schlussfolgerung}
\label{sec:conclusion}

Wir haben eine Umformulierung der Hodge-Vermutung durch die Perspektive von Positivitaet und No-Go-Obstruktionen vorgeschlagen. Das Schluesselkonzept -- \textit{arithmetische Positivitaet} -- kombiniert Hodge-Riemann-Positivitaet, absolute Stabilitaet unter $\mathrm{Aut}(\mathbb{C})$ und Lefschetz-Kompatibilitaet in eine einzige Bedingung.

Ein zentrales Ergebnis dieser Arbeit (Bemerkung~\ref{rem:ap_equals_absolute}) ist, dass der arithmetisch positive Kegel mit Delignes absoluten Hodge-Klassen uebereinstimmt: $\mathrm{AP}^p(X) = \mathrm{AbsHodge}^p(X)$. Die Arithmetische Positivitaetsvermutung ist daher aequivalent zu Delignes Frage ,,Sind alle absoluten Hodge-Klassen algebraisch?'' Die Positivitaetsbedingungen (AP1), (AP2b) und (AP3) werden fuer jede absolute Hodge-Klasse automatisch durch die Hodge-Riemann-Bilinearrelationen garantiert.

Obwohl die Arbeit keine genuein neue Vermutung einfuehrt, bietet sie:

\begin{enumerate}
\item \textbf{Eine neue Perspektive:} Die No-Go-Formulierung verwandelt die Hodge-Vermutung von einem Konstruktionsproblem in eine Unmoeglichkeitsaussage.

\item \textbf{Ein diagnostisches Werkzeug:} Das Galois-Hodge-Riemann-Spektrum macht explizit, welche Galois-Konjugation Nichtalgebraizitaet bezeugen koennte.

\item \textbf{Funktorialitaetsresultate:} Der arithmetisch positive Kegel ist funktoriell unter endlichem Rueckzug, endlichem Vorwaertsschub und K\"unneth-Produkt fuer primitive Klassen.
\end{enumerate}

Die verbleibende Herausforderung ist Delignes Frage: zu beweisen, dass absolute Hodge-Klassen algebraisch sind. Kuenftige Arbeit sollte staerkere Positivitaetsbedingungen suchen -- jenseits der Hodge-Riemann-Relationen -- die die Luecke zwischen absoluten Hodge-Klassen und algebraischen Klassen weiter einengen koennten.

\begin{quote}
\textit{,,Die Hodge-Vermutung verlangt nicht von uns, algebraische Zykel zu bauen. Sie verlangt von uns zu zeigen, dass die Geometrie keinen Raum fuer etwas anderes hat.''}
\end{quote}


% ===================================================================
% DANKSAGUNGEN
% ===================================================================
\begin{acknowledgments}
Diese Arbeit wurde als unabhaengige Forschung ohne institutionelle Foerderung durchgefuehrt.

\textit{KI-Werkzeuge.} KI-basierte Werkzeuge (Claude, Anthropic) wurden fuer mathematische Formalisierung und Texterstellung verwendet. Alle mathematischen Inhalte, Vermutungen und die Verantwortung fuer die Korrektheit liegen ausschliesslich beim Autor.

\textit{Foerderinformation.} Diese Forschung erhielt keine externe Foerderung.
\end{acknowledgments}


% ===================================================================
% BIBLIOGRAPHIE
% ===================================================================
\begin{thebibliography}{30}

\bibitem{Voisin2002}
C.~Voisin (2002).
\textit{Hodge Theory and Complex Algebraic Geometry I, II}.
Cambridge University Press.
DOI: 10.1017/CBO9780511615344.

\bibitem{GH1978}
P.~Griffiths, J.~Harris (1978).
\textit{Principles of Algebraic Geometry}.
John Wiley \& Sons, New York.

\bibitem{Deligne1982}
P.~Deligne (1982).
Hodge cycles on abelian varieties (Notes by J.~S.~Milne).
In \textit{Hodge Cycles, Motives, and Shimura Varieties}, Lecture Notes in Mathematics 900, pp.\ 9--100. Springer.
DOI: 10.1007/978-3-540-38955-2\_3.

\bibitem{Andre1996}
Y.~Andr\'e (1996).
Pour une th\'eorie inconditionnelle des motifs.
\textit{Publications Math\'ematiques de l'IH\'ES}, 83, 5--49.
DOI: 10.1007/BF02698643.

\bibitem{Charles2013}
F.~Charles (2013).
The Tate conjecture for K3 surfaces over finite fields.
\textit{Inventiones Mathematicae}, 194(1), 119--145.
DOI: 10.1007/s00222-012-0443-y.

\bibitem{Faltings1983}
G.~Faltings (1983).
Endlichkeitss\"atze f\"ur abelsche Variet\"aten \"uber Zahlk\"orpern.
\textit{Inventiones Mathematicae}, 73, 349--366.
DOI: 10.1007/BF01388432.

\bibitem{Abdulali1994}
S.~Abdulali (1994).
Algebraic cycles in families of abelian varieties.
\textit{Canadian Journal of Mathematics}, 46(6), 1121--1134.
DOI: 10.4153/CJM-1994-063-0.

\bibitem{Moonen1999}
B.~Moonen (1999).
Notes on Mumford-Tate groups.
Preprint, University of Amsterdam.

\bibitem{Voisin2002b}
C.~Voisin (2002).
A counterexample to the Hodge conjecture extended to K\"ahler varieties.
\textit{International Mathematics Research Notices}, 2002(20), 1057--1075.
DOI: 10.1155/S1073792802111135.

\bibitem{CattaniDeligneKaplan1995}
E.~Cattani, P.~Deligne \& S.~Kaplan (1995).
On the locus of Hodge classes.
\textit{Journal of the American Mathematical Society}, 8(2), 483--506.
DOI: 10.2307/2152824.

\bibitem{Grothendieck1969}
A.~Grothendieck (1969).
Standard conjectures on algebraic cycles.
In \textit{Algebraic Geometry (Bombay, 1968)}, pp.\ 193--199. Oxford University Press.

\bibitem{GrossZagier1986}
B.~Gross \& D.~Zagier (1986).
Heegner points and derivatives of $L$-series.
\textit{Inventiones Mathematicae}, 84, 225--320.
DOI: 10.1007/BF01388809.

\bibitem{Deligne1974}
P.~Deligne (1974).
La conjecture de Weil: I.
\textit{Publications Math\'ematiques de l'IH\'ES}, 43, 273--307.
DOI: 10.1007/BF02684373.

\bibitem{FST-Unified2026}
L.~Geiger (2026).
The Renormalized Free Energy Framework: A Unified Resolution of Structural Bottlenecks.
Zenodo preprint.
\href{https://doi.org/10.5281/zenodo.19036191}{doi:10.5281/zenodo.19036191}.

\bibitem{GeigerRH2026c}
L.~Geiger (2026).
The Riemann Hypothesis Part III: The Analytical Rank-Finite Certificate.
Zenodo preprint.
\href{https://doi.org/10.5281/zenodo.19035845}{doi:10.5281/zenodo.19035845}.

\bibitem{Mazur1972}
B.~Mazur (1972).
Frobenius and the Hodge filtration.
\textit{Bulletin of the American Mathematical Society}, 78(5), 653--667.
DOI: 10.1090/S0002-9904-1972-12976-8.

\bibitem{Moriwaki1999}
A.~Moriwaki (1999).
Arithmetic height functions over finitely generated fields.
\textit{Inventiones Mathematicae}, 140(1), 101--142.
DOI: 10.1007/s002220050358.

\bibitem{DeligneMilne1982}
P.~Deligne \& J.~S.~Milne (1982).
Tannakian categories.
In \textit{Hodge Cycles, Motives, and Shimura Varieties}, Lecture Notes in Mathematics 900, pp.\ 101--228. Springer.
DOI: 10.1007/978-3-540-38955-2\_4.

\bibitem{BhattScholze2022}
B.~Bhatt \& P.~Scholze (2022).
Prisms and prismatic cohomology.
\textit{Annals of Mathematics}, 196(3), 1135--1275.
DOI: 10.4007/annals.2022.196.3.5.

\bibitem{Bridgeland2007}
T.~Bridgeland (2007).
Stability conditions on triangulated categories.
\textit{Annals of Mathematics}, 166(2), 317--345.
DOI: 10.4007/annals.2007.166.317.

\bibitem{BBT2020}
B.~Bakker, B.~Brunebarbe \& J.~Tsimerman (2023).
o-minimal GAGA and a conjecture of Griffiths.
\textit{Inventiones Mathematicae}, 232(1), 163--228.
DOI: 10.1007/s00222-022-01166-1.

\end{thebibliography}

\end{document}
